\documentclass[12pt,a4paper]{article}
\usepackage[T1]{fontenc}
\usepackage[utf8]{inputenc}
\usepackage[german,italian]{babel}
\usepackage{amsmath,amssymb,amsthm}
\usepackage{geometry}
\usepackage{titlesec}
\usepackage{fancyhdr}
\usepackage{enumitem}
\usepackage{array}

\geometry{margin=1in}

\newcommand{\cO}{\mathcal{O}}
\newcommand{\cF}{\mathcal{F}}
\newcommand{\cE}{\mathcal{E}}
\newcommand{\cL}{\mathcal{L}}
\newcommand{\cH}{\mathcal{H}}
\newcommand{\cA}{\mathcal{A}}
\newcommand{\cB}{\mathcal{B}}
\newcommand{\cC}{\mathcal{C}}
\newcommand{\cD}{\mathcal{D}}
\newcommand{\cG}{\mathcal{G}}
\newcommand{\cI}{\mathcal{I}}
\newcommand{\cJ}{\mathcal{J}}
\newcommand{\cK}{\mathcal{K}}
\newcommand{\cM}{\mathcal{M}}
\newcommand{\cN}{\mathcal{N}}
\newcommand{\cP}{\mathcal{P}}
\newcommand{\cS}{\mathcal{S}}
\newcommand{\cT}{\mathcal{T}}
\newcommand{\cU}{\mathcal{U}}
\newcommand{\cV}{\mathcal{V}}
\newcommand{\cW}{\mathcal{W}}
\newcommand{\bA}{\mathbb{A}}
\newcommand{\bC}{\mathbb{C}}
\newcommand{\bF}{\mathbb{F}}
\newcommand{\bG}{\mathbb{G}}
\newcommand{\bN}{\mathbb{N}}
\newcommand{\bP}{\mathbb{P}}
\newcommand{\bQ}{\mathbb{Q}}
\newcommand{\bR}{\mathbb{R}}
\newcommand{\bZ}{\mathbb{Z}}
\newcommand{\fm}{\mathfrak{m}}
\newcommand{\fp}{\mathfrak{p}}
\newcommand{\fq}{\mathfrak{q}}
\newcommand{\fP}{\mathfrak{P}}
\newcommand{\fa}{\mathfrak{a}}
\DeclareMathOperator{\Gal}{Gal}
\DeclareMathOperator{\Aut}{Aut}
\DeclareMathOperator{\End}{End}
\DeclareMathOperator{\Hom}{Hom}
\DeclareMathOperator{\Ext}{Ext}
\DeclareMathOperator{\Tor}{Tor}
\DeclareMathOperator{\Pic}{Pic}
\DeclareMathOperator{\Ker}{Ker}
\let\Im\relax
\DeclareMathOperator{\Im}{Im}
\DeclareMathOperator{\codim}{codim}
\DeclareMathOperator{\supp}{supp}
\DeclareMathOperator{\Spec}{Spec}
\DeclareMathOperator{\Spf}{Spf}
\DeclareMathOperator{\Proj}{Proj}
\DeclareMathOperator{\Div}{Div}
\DeclareMathOperator{\ord}{ord}
\DeclareMathOperator{\card}{card}
\DeclareMathOperator{\rg}{rg}
\DeclareMathOperator{\tr}{tr}
\DeclareMathOperator{\Tr}{Tr}
\DeclareMathOperator{\Res}{Res}
\DeclareMathOperator{\Ind}{Ind}
\DeclareMathOperator{\Cor}{Cor}
\DeclareMathOperator{\Br}{Br}
\DeclareMathOperator{\Alb}{Alb}
\DeclareMathOperator{\Jac}{Jac}
\DeclareMathOperator{\Prym}{Prym}
\DeclareMathOperator{\Sym}{Sym}
\DeclareMathOperator{\Alt}{Alt}
\DeclareMathOperator{\Num}{Num}
\DeclareMathOperator{\GL}{GL}
\DeclareMathOperator{\tg}{tg}

\pagestyle{fancy}
\fancyhf{}
\fancyhead[C]{\small Riemann --- Gesammelte Werke (pp.~301--350)}
\fancyfoot[C]{\thepage}
\renewcommand{\headrulewidth}{0.4pt}

\titleformat{\section}{\large\bfseries}{\thesection.}{0.5em}{}
\titleformat{\subsection}{\normalsize\bfseries}{\S.~\thesubsection.}{0.5em}{}

\begin{document}

\begin{center}
\Large\bfseries Riemann --- Gesammelte Werke (pp.~301--350)
\end{center}

\tableofcontents
\newpage

%==============================================================================
% XIII. End of "Ueber die Hypothesen, welche der Geometrie zu Grunde liegen"
%==============================================================================
\section*{XIII. Ueber die Hypothesen, welche der Geometrie zu Grunde liegen.}
\addcontentsline{toc}{section}{XIII. Ueber die Hypothesen, welche der Geometrie zu Grunde liegen.}

\subsection*{3.}
\addcontentsline{toc}{subsection}{\S~3. In wie weit im Unendlichkleinen?}

\selectlanguage{german}

Die Fragen ueber das Unmessbargrosse sind fuer die Naturerklaerung muessige Fragen. Anders verhaelt es sich aber mit den Fragen ueber das Unmessbarkleine. Auf der Genauigkeit, mit welcher wir die Erscheinungen in's Unendlichkleine verfolgen, beruht wesentlich die Erkenntniss ihres Causalzusammenhangs. Die Fortschritte der letzten Jahrhunderte in der Erkenntniss der mechanischen Natur sind fast allein bedingt durch die Genauigkeit der Construction, welche durch die Erfindung der Analysis des Unendlichen und die von Archimed, Galil\"ai und Newton aufgefundenen einfachen Grundbegriffe, deren sich die heutige Physik bedient, m\"oglich geworden ist. In den Naturwissenschaften aber, wo die einfachen Grundbegriffe zu solchen Constructionen bis jetzt fehlen, verfolgt man, um den Causalzusammenhang zu erkennen, die Erscheinungen in's r\"aumlich Kleine, so weit es das Mikroskop nur gestattet. Die Fragen ueber die Massverh\"altnisse des Raumes im Unmessbarkleinen geh\"oren also nicht zu den muessigen.

Setzt man voraus, dass die K\"orper unabh\"angig vom Ort existiren, so ist das Kr\"ummungsmass ueberall constant, und es folgt dann aus den astronomischen Messungen, dass es nicht von Null verschieden sein kann; jedenfalls m\"usste sein reciprocer Werth eine Fl\"ache sein, gegen welche das unsern Teleskopen zug\"angliche Gebiet verschwinden m\"usste. Wenn aber eine solche Unabh\"angigkeit der K\"orper vom Ort nicht stattfindet, so kann man aus den Massverh\"altnissen im Grossen nicht auf die im Unendlichkleinen schliessen; es kann dann in jedem Punkte das Kr\"ummungsmass in drei Richtungen einen beliebigen Werth haben, wenn nur die ganze Kr\"ummung jedes messbaren Raumtheils nicht merklich von Null verschieden ist; noch complicirtere Verh\"altnisse k\"onnen eintreten, wenn die vorausgesetzte Darstellbarkeit eines Linienelements durch die Quadratwurzel aus einem Differentialausdruck zweiten Grades nicht stattfindet. Nun scheinen aber die empirischen Begriffe, in welchen die r\"aumlichen Massbestimmungen gegr\"undet sind, der Begriff des festen K\"orpers und des Lichtstrahls, im Unendlichkleinen ihre G\"ultigkeit zu verlieren; es ist also sehr wohl denkbar, dass die Massverh\"altnisse des Raumes im Unendlichkleinen den Voraussetzungen der Geometrie nicht gem\"ass sind, und dies w\"urde man in der That annehmen m\"ussen, sobald sich dadurch die Erscheinungen auf einfachere Weise erkl\"aren liessen.

Die Frage ueber die G\"ultigkeit der Voraussetzungen der Geometrie im Unendlichkleinen h\"angt zusammen mit der Frage nach dem Innern Grunde der Massverh\"altnisse des Raumes. Bei

\newpage
\noindent wohl noch zur Lehre vom R\"aume gerechnet werden darf, kommt die obige Bemerkung zur Anwendung, dass bei einer discreten Mannigfaltigkeit das Princip der Massverh\"altnisse schon in dem Begriffe dieser Mannigfaltigkeit enthalten ist, bei einer stetigen aber anders woher hinzukommen muss. Es muss also entweder das dem R\"aume zu Grunde liegende Wirkliche eine discrete Mannigfaltigkeit bilden, oder der Grund der Massverh\"altnisse ausserhalb, in darauf wirkenden bindenden Kr\"aften, gesucht werden.

Die Entscheidung dieser Fragen kann nur gefunden werden, indem man von der bisherigen durch die Erfahrung bew\"ahrten Auffassung der Erscheinungen, wozu Newton den Grund gelegt, ausgeht und diese durch Thatsachen, die sich aus ihr nicht erkl\"aren lassen, getrieben allm\"ahlich umarbeitet; solche Untersuchungen, welche, wie die hier gef\"uhrte, von allgemeinen Begriffen ausgehen, k\"onnen nur dazu dienen, dass diese Arbeit nicht durch die Beschr\"anktheit der Begriffe gehindert und der Fortschritt im Erkennen des Zusammenhangs der Dinge nicht durch ueberlieferte Vorurtheile gehemmt wird.

Es f\"uhrt dies hin\"uber in das Gebiet einer andern Wissenschaft, in das Gebiet der Physik, welches wohl die Natur der heutigen Veranlassung nicht zu betreten erlaubt.

\medskip
\noindent\textbf{"Ubersicht.}

\begin{center}
\begin{tabular}{p{10cm}r}
\multicolumn{2}{r}{Seite} \\[2mm]
Plan der Untersuchung & 272 \\[1mm]
I. Begriff einer $n$fach ausgedehnten Gr\"osse\textsuperscript{*}) & 273 \\[1mm]
\quad \S.~1. Stetige und discrete Mannigfaltigkeiten. Bestimmte Theile einer Mannigfaltigkeit heissen Quanta. Eintheilung der Lehre von den stetigen Gr\"ossen in die Lehre & \\[1mm]
\quad 1) von den blossen Gebietsverh\"altnissen, bei welcher eine Unabh\"angigkeit der Gr\"ossen vom Ort nicht vorausgesetzt wird, & \\[1mm]
\quad 2) von den Massverh\"altnissen, bei welcher eine solche Unabh\"angigkeit vorausgesetzt werden muss & 273 \\[2mm]
\quad \S.~2. Erzeugung des Begriffs einer einfach, zweifach, \dots, $n$fach ausgedehnten Mannigfaltigkeit & 275 \\[1mm]
\quad \S.~3. Zur\"uckf\"uhrung der Ortsbestimmung in einer gegebenen Mannigfaltigkeit auf Quantit\"atsbestimmungen. Wesentliches Kennzeichen einer $n$fach ausgedehnten Mannigfaltigkeit & 275 \\[2mm]
II. Massverh\"altnisse, deren eine Mannigfaltigkeit von $n$ Dimensionen f\"ahig
\end{tabular}
\end{center}

\newpage
\begin{center}
\begin{tabular}{p{10cm}r}
\noindent ist\textsuperscript{*}), unter der Voraussetzung, dass die Linien unabh\"angig von der Lage eine L\"ange besitzen, also jede Linie durch jede messbar ist & 276 \\[2mm]
\quad \S.~1. Ausdruck des Linienelements. Als eben werden solche Mannigfaltigkeiten betrachtet, in denen das Linienelement durch die Wurzel aus einer Quadratsumme vollst\"andiger Differentialien ausdr\"uckbar ist & 277 \\[2mm]
\quad \S.~2. Untersuchung der $n$fach ausgedehnten Mannigfaltigkeiten, in welchen das Linienelement durch die Quadratwurzel aus einem Differentialausdruck zweiten Grades dargestellt werden kann. Mass ihrer Abweichung von der Ebenheit (Kr\"ummungsmass) in einem gegebenen Punkte und einer gegebenen Fl\"achenrichtung. Zur Bestimmung ihrer Massverh\"altnisse ist es (unter gewissen Beschr\"ankungen) zul\"assig und hinreichend, dass das Kr\"ummungsmass in jedem Punkte in $\frac{1}{2}n(n-1)$ Fl\"achenrichtungen beliebig gegeben wird & 279 \\[2mm]
\quad \S.~3. Geometrische Erl\"auterung & 280 \\[1mm]
\quad \S.~4. Die ebenen Mannigfaltigkeiten (in denen das Kr\"ummungsmass allenthalben $=0$ ist) lassen sich betrachten als einen besondern Fall der Mannigfaltigkeiten mit constantem Kr\"ummungsmass. Diese k\"onnen auch dadurch definirt werden, dass in ihnen Unabh\"angigkeit der $n$fach ausgedehnten Gr\"ossen vom Ort (Bewegbarkeit derselben ohne Dehnung) stattfindet & 281 \\[2mm]
\quad \S.~5. Fl\"achen mit constantem Kr\"ummungsmasse & 282 \\[2mm]
III. Anwendung auf den Raum & 283 \\[2mm]
\quad \S.~1. Systeme von Thatsachen, welche zur Bestimmung der Massverh\"altnisse des Raumes, wie die Geometrie sie voraussetzt, hinreichen & 283 \\[1mm]
\quad \S.~2. In wie weit ist die G\"ultigkeit dieser empirischen Bestimmungen wahrscheinlich jenseits der Grenzen der Beobachtung im Unmessbargrossen? & 283 \\[1mm]
\quad \S.~3. In wie weit im Unendlichkleinen? Zusammenhang dieser Frage mit der Naturerklaerung\textsuperscript{**}) & 286
\end{tabular}
\end{center}

\bigskip
\noindent\small\textsuperscript{*}) Art.~I. bildet zugleich die Vorarbeit f\"ur Beitr\"age zur analysis situs.

\small\textsuperscript{**}) Der \S.~3 des Art.~III bedarf noch einer Umarbeitung und weiteren Ausf\"uhrung. (Bemerkungen von Riemann.)

\newpage
%==============================================================================
% XIV. Ein Beitrag zur Elektrodynamik
%==============================================================================
\section*{XIV. Ein Beitrag zur Elektrodynamik.}
\addcontentsline{toc}{section}{XIV. Ein Beitrag zur Elektrodynamik.}

\noindent\small{(Aus Poggendorff's Annalen der Physik und Chemie, Bd.~CXXXI.)}

\medskip

\normalsize
Der K\"oniglichen Societ\"at erlaube ich mir eine Bemerkung mitzutheilen, welche die Theorie der Elektricit\"at und des Magnetismus mit der des Lichts und der strahlenden W\"arme in einen nahen Zusammenhang bringt. Ich habe gefunden, dass die elektrodynamischen Wirkungen galvanischer Str\"ome sich erkl\"aren lassen, wenn man annimmt, dass die Wirkung einer elektrischen Masse auf die uebrigen nicht momentan geschieht, sondern sich mit einer constanten (der Lichtgeschwindigkeit innerhalb der Grenzen der Beobachtungsfehler gleichen) Geschwindigkeit zu ihnen fortpflanzt. Die Differentialgleichung f\"ur die Fortpflanzung der elektrischen Kraft wird bei dieser Annahme dieselbe, wie die f\"ur die Fortpflanzung des Lichts und der strahlenden W\"arme.

Es seien $S$ und $S'$ zwei von constanten galvanischen Str\"omen durchflossene und gegen einander nicht bewegte Leiter, $s$ sei ein elektrisches Massentheilchen im Leiter $S$, welches sich zur Zeit $t$ im Punkte $(x,y,z)$ befinde, $e$ ein elektrisches Massentheilchen von $S'$ und befinde sich zur Zeit $t$ im Punkte $(x',y',z')$. Ueber die Bewegung der elektrischen Massentheilchen, welche in jedem Leitertheilchen f\"ur die positiv und negativ elektrischen entgegengesetzt ist, mache ich die Voraussetzung, dass sie in jedem Augenblicke so vertheilt sind, dass die Summen
\[\sum s f(x,y,z), \quad \sum s' f(x',y',z'),\]
\"uber s\"ammtliche Massentheilchen der Leiter ausgedehnt gegen dieselben Summen, wenn sie nur ueber die positiv elektrischen oder nur ueber die negativ elektrischen Massentheilchen ausgedehnt werden, vernachl\"assigt werden d\"urfen, sobald die Function $f$ und ihre Differentialquotienten stetig sind.

Diese Voraussetzung kann auf sehr mannigfaltige Weise erf\"ullt werden. Nimmt man z.~B. an, dass die Leiter in den kleinsten Theilen krystallinisch sind, so dass sich dieselbe relative Vertheilung der

\newpage
\noindent Elektricit\"aten in bestimmten gegen die Dimensionen der Leiter unendlich kleinen Abst\"anden periodisch wiederholt, so sind, wenn $\beta$ die L\"ange einer solchen Periode bezeichnet, jene Summen unendlich klein, wie $\beta^{n}$, wenn $f$ und ihre Derivirten bis zur $(n-1)$ten Ordnung stetig sind, und unendlich klein wie $e^{-\frac{c}{\beta}}$, wenn sie s\"ammtlich stetig sind.

\medskip
\noindent\textbf{Erfahrungsm\"assiges Gesetz der elektrodynamischen Wirkungen.}

Sind die specifischen Stromintensit\"aten nach mechanischem Mass zur Zeit $t$ im Punkte $(x,y,z)$ parallel den drei Axen $u,v,w$, und im Punkte $(x',y',z')$ $u',v',w'$, und bezeichnet $r$ die Entfernung beider Punkte, $c$ die von Kohlrausch und Weber bestimmte Constante, so ist der Erfahrung nach das Potential der von $S$ auf $S'$ ausge\"ubten Kr\"afte
\[\frac{2}{cc} \int\!\int \frac{uu' + vv' + ww'}{r}\, dS\, dS',\]
dieses Integral ueber s\"ammtliche Elemente $dS$ und $dS'$ der Leiter $S$ und $S'$ ausgedehnt. F\"uhrt man statt der specifischen Stromintensit\"aten die Producte aus den Geschwindigkeiten in die specifischen Dichtigkeiten und dann f\"ur die Producte aus diesen in die Volumelemente die in ihnen enthaltenen Massen ein, so geht dieser Ausdruck ueber in
\[\frac{1}{cc} \sum\sum \frac{ee'}{r} \frac{d d'(r^{2})}{dt\,dt'},\]
wenn die Aenderung von $r^{2}$ w\"ahrend der Zeit $dt$, welche von der Bewegung von $e$ herr\"uhrt, durch $d$, und die von der Bewegung von $e'$ herr\"uhrende durch $d'$ bezeichnet wird.

Dieser Ausdruck kann durch Hinwegnahme von
\[\sum e' \sum \frac{s}{cc} \frac{d'(r^{2})}{r} \frac{1}{dt}\, dt',\]
welches durch die Summirung nach $e$ verschwindet, in
\[\frac{1}{cc} \sum\sum \frac{e'e}{r} \frac{d'(r^{2})}{dt}\, dt\]
und dieses wieder durch Addition von
\[\sum e \sum \frac{e'}{cc} \frac{d'(r^{2})}{r} \frac{1}{dt}\, dt',\]
welches durch die Summation nach $e'$ Null wird, in
\[\frac{1}{cc} \sum\sum ee' \frac{dd'(r^{2})}{dt\, dt'}\]
verwandelt werden.

\medskip
\noindent\textbf{Ableitung dieses Gesetzes aus der neuen Theorie.}

Nach der bisherigen Annahme ueber die elektrostatische Wirkung wird die Potentialfunction $U$ beliebig vertheilter elektrischer Massen, wenn $\varrho$ ihre Dichtigkeit im Punkte $(x,y,z)$ bezeichnet, durch die Bedingung
\[\frac{\partial^{2}U}{\partial x^{2}} + \frac{\partial^{2}U}{\partial y^{2}} + \frac{\partial^{2}U}{\partial z^{2}} + 4\pi\varrho = 0\]
und durch die Bedingung, dass $U$ stetig und in unendlicher Entfernung von wirkenden Massen constant sei, bestimmt. Ein particul\"ares Integral der Gleichung
\[\frac{\partial^{2}U}{\partial x^{2}} + \frac{\partial^{2}U}{\partial y^{2}} + \frac{\partial^{2}U}{\partial z^{2}} = 0\]
welches ueberall ausser dem Punkte $(x',y',z')$ stetig bleibt, ist
\[\frac{1}{r},\]
und diese Function bildet die vom Punkte $(x',y',z')$ aus erzeugte Potentialfunction, wenn sich in demselben zur Zeit $t$ die Masse $-f(t)$ befindet.

Statt dessen nehme ich nun an, dass die Potentialfunction $U$ durch die Bedingung
\[\frac{\partial^{2}U}{\partial x^{2}} + \frac{\partial^{2}U}{\partial y^{2}} + \frac{\partial^{2}U}{\partial z^{2}} - aa \frac{\partial^{2}U}{\partial t^{2}} + 4\pi\varrho = 0\]
bestimmt wird, so dass die vom Punkte $(x',y',z')$ aus erzeugte Potentialfunction, wenn sich in demselben zur Zeit $t$ die Masse $-f(t)$ befindet,
\[\frac{f\left(t - \frac{r}{a}\right)}{r}\]
wird.

Bezeichnet man die Coordinaten der Masse $e$ zur Zeit $t$ durch $x_{t}$, $y_{t}$, $z_{t}$, und die der Masse $e'$ zur Zeit $t'$ durch $x'_{t'}$, $y'_{t'}$, $z'_{t'}$, und setzt zur Abk\"urzung
\[\frac{1}{a}\sqrt{(x_{t} - x'_{t'})^{2} + (y_{t} - y'_{t'})^{2} + (z_{t} - z'_{t'})^{2}} = \frac{r}{a} = F(t,t'),\]
so wird nach dieser Annahme das Potential von $e$ auf $e'$ zur Zeit $t$
\[= -ee' \frac{f'(t - aF(t,t'))}{r}.\]

\newpage
Das Potential der von s\"ammtlichen Massen $e$ des Leiters $S$ auf die Massen $e'$ des Leiters $S'$ von der Zeit $0$ bis zur Zeit $t$ ausge\"ubten Kr\"afte wird daher
\[P = - \int_{0}^{t} \sum\sum ee' F'(t-a, t')\, dt',\]
die Summen ueber s\"ammtliche Massen beider Leiter ausgedehnt.

Da die Bewegung f\"ur entgegengesetzt elektrische Massen in jedem Leitertheilchen entgegengesetzt ist, so erlangt die Function $F(t,t')$ durch die Derivation nach $t$ die Eigenschaft, mit $e$, und durch die Derivation nach $t'$ die Eigenschaft, mit $e'$ ihr Zeichen zu \"andern. Bei der vorausgesetzten Vertheilung der Elektricit\"aten wird daher, wenn man die Derivationen nach $t$ durch obere und nach $t'$ durch untere Accente bezeichnet, $\sum\sum ee' F'_{n}(t,t)$, ueber s\"ammtliche elektrische Massen ausgedehnt, nur dann nicht unendlich klein gegen die ueber die elektrischen Massen einer Art erstreckte Summe, wenn $n$ und $\nu$ beide ungerade sind.

Man nehme nun an, dass die elektrischen Massen w\"ahrend der Fortpflanzungszeit der Kraft von einem Leiter zum anderen nur einen sehr kleinen Weg zur\"ucklegen, und betrachte die Wirkung w\"ahrend eines Zeitraums, gegen welchen die Fortpflanzungszeit verschwindet. In dem Ausdrucke von $P$ kann man dann zun\"achst
\[F'(t-a, t') - F'(t, t')\]
durch
\[-a \frac{\partial F'(t,t')}{\partial t}\]
ersetzen, da $\sum\sum ee' F'(t,t)$ vernachl\"assigt werden darf. Man erh\"alt dadurch
\[P = \int_{0}^{t} dt' \sum\sum ee' a \frac{\partial F'(t,t')}{\partial t},\]
oder wenn man die Ordnung der Integrationen umkehrt und $t+a$ f\"ur $\tau$ setzt,
\[P = \sum\sum ee' \int_{0}^{t} d\vartheta \int_{-\vartheta}^{t} d\tau\, F'(\tau, \tau+\vartheta).\]
Verwandelt man die Grenzen des innern Integrals in $0$ und $t$, so wird dadurch an der obern Grenze der Ausdruck

\newpage
\noindent $H(t) = \sum\sum ee' \int_{0}^{t} d\vartheta \int_{-\vartheta}^{0} dt\, F'(t+\tau, t+\tau+\vartheta)$

\medskip
\noindent hinzugef\"ugt, und an der untern Grenze der Werth dieses Ausdrucks f\"ur $t=0$ hinweggenommen. Man hat also
\[P = \int_{0}^{t} dt' \sum\sum ee' \int_{0}^{t} d\vartheta\, F'(t, \tau+\vartheta) - H(t) + H(0).\]
In diesem Ausdrucke kann man $F'(t, t+\vartheta)$ durch $F'(t, t+\vartheta) - F'(\tau, \tau)$ ersetzen, da
\[\sum\sum ee' F'(\tau, \tau)\]
vernachl\"assigt werden darf. Man erh\"alt dadurch als Factor von $ee'$ einen Ausdruck, der sowohl mit $e$ als mit $e'$ sein Zeichen \"andert, so dass sich bei den Summationen die Glieder nicht gegen einander aufheben, und unendlich kleine Bruchtheile der einzelnen Glieder vernachl\"assigt werden d\"urfen. Es ergiebt sich daher, indem man
\[F'(t, \tau+\vartheta) - F'(t, \tau)\]
durch
\[\vartheta \frac{\partial F'(t,\tau)}{\partial \tau}\]
ersetzt und die Integration nach $\vartheta$ ausf\"uhrt, bis auf einen zu vernachl\"assigenden Bruchtheil
\[P = \int_{0}^{t} \frac{aa}{2} \sum\sum ee' \frac{\partial F'(t,\tau)}{\partial \tau}\, dt.\]
Es ist leicht zu sehen, dass $H(t)$ und $H(0)$ vernachl\"assigt werden d\"urfen; denn es ist
\[H(t) = \sum\sum ee' \int_{0}^{t} d\vartheta \int_{-\vartheta}^{0} d\tau\, F'(t+\tau, t+\tau+\vartheta),\]
folglich:
\[H'(t) = \sum\sum ee' \int_{0}^{t} d\vartheta \left(F'(t,t+\vartheta) - F'(t-\vartheta, t)\right).\]
Hierin aber ist nur das erste Glied des Factors von $ee'$ mit dem Factor in dem ersten Bestandtheile von $P$ von gleicher Ordnung, und dieses liefert wegen der Summation nach $e$ nur einen zu vernachl\"assigenden Bruchtheil desselben.

Der Werth von $P$, welcher sich aus unserer Theorie ergiebt, stimmt mit dem erfahrungsm\"assigen
\[\frac{1}{cc} \sum\sum ee' \frac{dd'(r^{2})}{dt\,dt'}\]
\"uberein, wenn man $aa = \frac{1}{4}cc$ annimmt.

Nach der Bestimmung von Weber und Kohlrausch ist
\[c = 439450 \cdot 10^{8} \frac{\text{mm}}{\text{Sec}},\]
woraus sich $a$ zu $41949$ geographischen Meilen in der Secunde ergiebt, w\"ahrend f\"ur die Lichtgeschwindigkeit von Busch aus Bradley's Aberrationsbeobachtungen $41994$ Meilen, und von Fizeau durch directe Messung $41882$ Meilen gefunden worden sind.

\medskip
\noindent\small Dieser Aufsatz wurde von Riemann der K\"onigl.\ Gesellschaft der Wissenschaften zu G\"ottingen am 10.\ Februar 1858 \"uberreicht, wie aus einer dem Titel des Manuscriptes hinzugef\"ugten Bemerkung des damaligen Secret\"ars der Gesellschaft hervorgeht, sp\"ater aber wieder zur\"uckgezogen. Nachdem der Aufsatz nach Riemann's Tode ver\"offentlicht worden war, wurde er durch Clausius (Poggendorff's Annalen Bd.~CXXXV p.~606) einer Kritik unterworfen, deren wesentlichster Einwand in Folgendem besteht:

Nach den Voraussetzungen hat die Summe:
\[\sum\sum \frac{ee'}{r}\]
einen verschwindend kleinen Werth. Die Operation, verm\"oge deren sp\"ater daf\"ur ein nicht verschwindend kleiner Werth gefunden wird, muss daher einen Irrthum enthalten, den Clausius in der Ausf\"uhrung einer unberechtigten Umkehrung der Integrationsfolge findet.

Der Einwand scheint mir begr\"undet und ich bin mit Clausius der Meinung, dass Riemann sich denselben selbst gemacht und deshalb die Arbeit vor der Publication zur\"uckgezogen hat.

Obwohl damit der wesentlichste Inhalt der Riemann'schen Deduction dahinfallen w\"urde, habe ich mich doch zur Aufnahme dieses Aufsatzes in die vorliegende Sammlung entschlossen, weil ich nicht zu entscheiden wagte, ob er nicht doch noch Keime zu weiteren fruchtbaren Gedanken ueber diese h\"ochst interessante Frage enth\"alt. W.

\newpage
%==============================================================================
% XV. Beweis des Satzes
%==============================================================================
\section*{XV. Beweis des Satzes, dass eine einwerthige mehr als $2n$fach periodische Function von $n$ Ver\"anderlichen unm\"oglich ist.}
\addcontentsline{toc}{section}{XV. Beweis des Satzes, dass eine einwerthige mehr als $2n$fach periodische Function unm\"oglich ist.}

\noindent\small{(Auszug aus einem Schreiben Riemann's an Herrn Weierstrass.)}

\noindent\small{(Aus Borchardt's Journal f\"ur reine und angewandte Mathematik, Bd.~71.)}

\medskip

\normalsize
\dots\ Den Beweis des Satzes, auf welchen Sie neulich die Unterhaltung lenkten, dass eine einwerthige mehr als $2n$fach periodische Function von $n$ Ver\"anderlichen unm\"oglich ist, habe ich im Gespr\"ach wohl nicht ganz klar ausgedr\"uckt, auch nur die Grundgedanken angegeben; ich theile ihn Ihnen daher hier noch einmal mit.

Es sei $f$ eine $2n$fach periodische Function von $n$ Ver\"anderlichen $x_{1}, x_{2}, \dots, x_{n}$ und --- ich darf wohl meine Ihnen bekannten Benennungen gebrauchen --- der Periodicit\"atsmodul von $x_{\nu}$ f\"ur die $\mu$te Periode $a^{\mu}_{\nu}$. Es lassen sich dann bekanntlich die Gr\"ossen $x$ in die Form
\[x_{\nu} = \sum_{\mu} a^{\mu}_{\nu} \xi_{\mu} \quad \text{f\"ur } \nu = 1, 2, \dots, n\]
setzen\textsuperscript{*}), so dass die Gr\"ossen $\xi$ reell sind. L\"asst man nun die Gr\"ossen $\xi$ die Werthe von $0$ bis $1$ mit Ausschluss eines von diesen Grenzwerthen durchlaufen, so hat das dadurch entstehende $2n$fach ausgedehnte Gr\"ossengebiet die Eigenschaft, dass jedes System von Werthen der $n$ Ver\"anderlichen einem und nur einem Werthsysteme innerhalb dieses Gr\"ossengebiets nach den $2n$ Mod\"ulsystemen congruent ist. Ich werde, um mich sp\"ater k\"urzer ausdr\"ucken zu k\"onnen, dieses Gebiet ``das bei diesen $2n$ Modulsystemen periodisch sich wiederholende Gr\"ossengebiet'' nennen.

Hat die Function nun noch ein $(2n+1)$tes Modulsystem, welches sich nicht aus den $2n$ ersten Modulsystemen zusammensetzen l\"asst, so kann man die einem Gr\"ossensysteme nach diesem Modulsysteme congruenten Gr\"ossensysteme auf innerhalb dieses Gebiets liegende nach den $2n$ ersten Modulsystemen ihnen congruente zur\"uckf\"uhren und dadurch offenbar beliebig viele innerhalb dieses Gebiets liegende und nach den $2n+1$ Modulsystemen einander congruente Gr\"ossensysteme erhalten, wenn nicht zwei von den nach dem $(2n+1)$ten Modulsysteme congruente Gr\"ossensysteme auch nach den $2n$ ersten Modulsystemen congruent sind. In diesem Falle w\"urden zwischen den $2n+1$ Modulsystemen $n$ Gleichungen von der Form
\[\xi_{\nu} = \sum_{\mu=1}^{2n+1} m^{\mu} a^{\mu}_{\nu}\]
worin die Gr\"ossen $m$ ganze Zahlen w\"aren, stattfinden, und folglich, wie ich sp\"ater zeigen werde, die $2n+1$ Modulsysteme sich aus $2n$ Modulsystemen zusammensetzen lassen.

Man theile nun f\"ur jede der Gr\"ossen $\xi$ die Strecke von $0$ bis $1$ in $q$ gleiche Theile, wodurch das bei den $2n$ ersten Modulsystemen periodisch wiederkehrende Gebiet in $q^{2n}$ Gebiete zerf\"allt, in deren jedem sich die Gr\"ossen $\xi$ nur um $\frac{1}{q}$ \"andern. Offenbar m\"ussen dann von mehr als $q^{2n}$ nach den $2n+1$ Modulsystemen einander congruenten und in jenem Gebiete liegenden Gr\"ossensystemen nothwendig zwei in dasselbe Theilgebiet fallen, so dass sich die Werthe derselben Gr\"osse $\xi$ in beiden keinesfalls um mehr als $\frac{1}{q}$ von einander unterscheiden. Die Function bleibt also dann unge\"andert, w\"ahrend keine der Gr\"ossen $\xi$ um mehr als $\frac{1}{q}$ ge\"andert wird, und ist folglich, da $q$ beliebig gross genommen werden kann, wenn sie stetig ist, eine Function von weniger als $n$ linearen Ausdr\"ucken der Gr\"ossen $x$.

Es ist nun noch zu zeigen, dass sich $2n+1$ Modulsysteme, zwischen denen die $n$ Gleichungen
\[\xi_{\nu} = \sum_{\mu=1}^{2n+1} m^{\mu} a^{\mu}_{\nu}\]
stattfinden, aus $2n$ Modulsystemen zusammensetzen lassen.

Man kann zun\"achst leicht beweisen, dass sich zu einem Modulsysteme
\[\frac{a_{\nu 1}}{\vartheta_{1}}, \frac{a_{\nu 2}}{\vartheta_{2}}, \dots, \frac{a_{\nu, 2n}}{\vartheta_{2n}},\]
worin die Gr\"ossen $m$ ganze Zahlen ohne gemeinschaftlichen Theiler sind, immer $2n-1$ andere Modulsysteme $b_{1}, b_{2}, \dots, b_{2n}$ so finden lassen,

\newpage
\noindent dass Congruenz nach den Modulsystemen $a$ mit Congruenz nach den Modulsystemen $b$ identisch ist. Es seien $\vartheta_{1}$ der gr\"osste gemeinschaftliche Theiler von $m_{1}$ und $m_{2}$ und $\alpha$, $\beta$ zwei der Gleichung
\[\beta m_{1} - \alpha m_{2} = \vartheta_{1}\]
gen\"ugende ganze Zahlen. Setzt man dann
\[\alpha a_{\nu 1} + \beta a_{\nu 2} = c_{\nu 1}\]
und
\[\alpha a'_{\nu 1} + \beta a'_{\nu 2} = b_{\nu, 2n},\]
so hat man
\[a_{\nu 1} = \beta c_{\nu 1} - \frac{m_{2}}{\vartheta_{1}} b_{\nu, 2n}, \quad a_{\nu 2} = -\alpha c_{\nu 1} + \frac{m_{1}}{\vartheta_{1}} b_{\nu, 2n}.\]
Es lassen sich also auch umgekehrt die Modulsysteme $a_{1}$ und $a_{2}$ aus den Modulsystemen $b_{\nu, 2n}$ und $c_{\nu 1}$ zusammensetzen, und folglich ist Congruenz nach jenen mit Congruenz nach diesen gleichbedeutend. Man kann daher die Modulsysteme $a_{1}$ und $a_{2}$ durch die Modulsysteme $c_{1}$ und $b_{2n}$ ersetzen. Auf dieselbe Weise kann man nun, wenn $\vartheta_{2}$ der gr\"osste gemeinschaftliche Theiler von $\vartheta_{1}$ und $m_{3}$ ist, die Modulsysteme $c_{1}$ und $a_{3}$ durch das Modulsystem
\[\frac{c_{\nu 1}}{\vartheta_{2}}\]
und durch ein Modulsystem $b_{\nu, 2n-1}$ ersetzen. Durch Fortsetzung dieses Verfahrens erh\"alt man offenbar den zu beweisenden Satz. Der Inhalt des periodisch sich wiederholenden Gebiets ist f\"ur die neuen Modulsysteme $b$ derselbe wie f\"ur die alten.

Mit H\"ulfe dieses Satzes lassen sich in den $n$ Gleichungen
\[\xi_{\nu} = \sum_{\mu=1}^{2n+1} m^{\mu} a^{\mu}_{\nu}\]
die $2n$ ersten Modulsysteme so durch $2n$ neue $b_{1}, b_{2}, \dots, b_{2n}$ ersetzen, dass diese Gleichungen die Form
\[\xi_{\nu} = \frac{p\, b^{\,2n+1}_{\nu} + q\, b^{\,1}_{\nu}}{\vartheta}\]
annehmen, worin $p$ und $q$ ganze Zahlen ohne gemeinschaftlichen Theiler sind. Sind nun $\gamma$, $\delta$ zwei der Gleichung
\[p\delta + q\gamma = 1\]
gen\"ugende ganze Zahlen, so lassen sich offenbar die beiden Modulsysteme $b^{1}$ und $a^{2n+1}$ durch das eine Modulsystem
\[\gamma\, b^{\,2n+1}_{\nu} + \delta\, b^{\,1}_{\nu}\]
ersetzen. S\"ammtliche Modulsysteme, welche sich aus den Modulsystemen $a_{1}, a_{2}, \dots, a_{2n+1}$ zusammensetzen lassen, k\"onnen also auch aus den $2n$ Modulsystemen $\frac{b}{\vartheta}$, $b_{2}, b_{3}, \dots, b_{2n}$ zusammengesetzt werden, und umgekehrt. Der Inhalt des periodisch wiederkehrenden Gebiets betr\"agt f\"ur diese $2n$ Modulsysteme nur $\frac{1}{\vartheta}$ von dem f\"ur die $2n$ ersten Modulsysteme $a$. Hat die Function nun ausser diesen Modulsystemen noch ein durch \"ahnliche ganzzahlige Gleichungen mit ihnen verbundenes, so lassen sich wieder $2n$ neue Modulsysteme finden, aus welchen sich alle diese Modulsysteme zusammensetzen lassen, und der Inhalt des periodisch sich wiederholenden Gebiets wird dabei wieder auf einen aliquoten Theil reducirt. Wenn dieses Gebiet unendlich klein wird, so wird die Function eine Function von weniger als $n$ linearen Ausdr\"ucken der Ver\"anderlichen und zwar von $n-1$ oder $n-2$ oder $n-m$, jenachdem nur eine, oder zwei oder $m$ Dimensionen dieses Gr\"ossengebiets unendlich klein werden. Soll dies aber nicht eintreten, so muss die Operation schliesslich abbrechen, und man wird also zu $2n$ Modulsystemen gelangen, aus welchen sich s\"ammtliche Modulsysteme der Function zusammensetzen lassen.

\bigskip
\noindent G\"ottingen, den 26ten October 1859.

\bigskip
\noindent\small\textsuperscript{*}) Dies ist nicht immer der Fall, sondern nur, wenn die $2n$ Gleichungen, durch welche die Gr\"ossen $\xi$ bestimmt werden, von einander unabh\"angig sind; die Ausnahmen sind aber leicht zu behandeln.

\newpage
%==============================================================================
% XVI. Extract from Italian letter to Betti
%==============================================================================
\section*{XVI. Estratto di una lettera scritta in lingua Italiana il di 21 Gennaio 1864 al Sig.\ Professore Enrico Betti.}
\addcontentsline{toc}{section}{XVI. Estratto di una lettera al Sig.\ Professore Enrico Betti.}

\noindent\small{(Annali di Matematica, Ser.~1, T.~VII.)}

\medskip

\selectlanguage{italian}

Carissimo Amico

\dots Per trovare l'attrazione di un cilindro omogeneo retto ellisoidale qualunque, io considero, introducendo coordinate rettangolari $x$, $y$, $z$, il cilindro infinito limitato della diseguaglianza:
\[\frac{x^{2}}{a^{2}} + \frac{y^{2}}{b^{2}} + \frac{z^{2}}{c^{2}} < 1,\]
ripieno di massa di densit\`a costante $+1$, se $\frac{x^{2}}{a^{2}} + \frac{y^{2}}{b^{2}} + \frac{z^{2}}{c^{2}} < 0$, e di densit\`a $-1$, se $\frac{x^{2}}{a^{2}} + \frac{y^{2}}{b^{2}} + \frac{z^{2}}{c^{2}} > 0$. Allora se poniamo, come \`e solito, il potenziale nel punto $x$, $y$, $z$ eguale a $F$ e
\[\frac{\partial F}{\partial x} = X, \quad \frac{\partial F}{\partial y} = Y, \quad \frac{\partial F}{\partial z} = Z,\]
si ha per $\frac{x^{2}}{a^{2}} + \frac{y^{2}}{b^{2}} < 0$, $F = 0$, $X = 0$, $Y = 0$.

$Z$ \`e eguale al potenziale dell' ellisse:
\[1 - \frac{x^{2}}{a^{2}} - \frac{y^{2}}{b^{2}} > 0\]
colla densit\`a $2$, e si trova col metodo di Dirichlet, se denotiamo con $\mathfrak{C}$ la radice maggiore dell' equazione:
\[\frac{x^{2}}{a^{2}+s} + \frac{y^{2}}{b^{2}+s} + \frac{z^{2}}{c^{2}+s} = 1,\]
\[Z = 2abc \int_{-\mathfrak{C}}^{\infty} \frac{\sqrt{(a^{2}+s)(b^{2}+s)}}{(c^{2}+s)^{2}} \, \frac{ds}{D},\]
con $D$:
\[D = \sqrt{(a^{2}+s)(b^{2}+s)(c^{2}+s)}.\]

$X$ ed $Y$ si possono determinare dalle equazioni:
\[\frac{\partial X}{\partial z} = \frac{\partial Z}{\partial x}, \quad \frac{\partial Y}{\partial z} = \frac{\partial Z}{\partial y}\]
e dalle condizioni:
\[X = 0, \quad Y = 0\]
per $z = 0$.

Per effettuare questa determinazione conviene di sostituire invece di $Z$ la sua espressione esteso per il contorno intero di un pezzo del Piano degli $s$, che contiene il valore $\mathfrak{C}$ senza contenere verun altro valore di diramazione o di discontinuit\`a della funzione sotto il segno integrale. Se denotiamo le radici di $F=0$ in ordine di grandezza con $\mathfrak{a}$, $\mathfrak{a}'$, $\mathfrak{a}''$ questi valori sono tutti reali e in ordine di grandezza:
\[\mathfrak{a}, \; \mathfrak{a}', \; -b^{2}, \; \mathfrak{a}'', \; -a^{2},\]
in modo che:
\[\mathfrak{a} > 0 > \mathfrak{a}' > -b^{2} > \mathfrak{a}'' > -a^{2}.\]

Posto
\[F = \mathfrak{a}' - \frac{1}{s},\]
viene
\[Z = 2abc \int_{0}^{\mathfrak{a}'} \frac{\sqrt{(a^{2}+s)(b^{2}+s)}}{(c^{2}+s)^{2}} \, \frac{ds}{D\sqrt{F}},\]
\[X = \int \frac{\partial Z}{\partial x}\, dz, \quad Y = \int \frac{\partial Z}{\partial y}\, dz.\]

Ma:
\[\frac{\partial Z}{\partial x} = \frac{\partial Z}{\partial t} \frac{\partial t}{\partial x}, \quad \frac{\partial Z}{\partial y} = \frac{\partial Z}{\partial t} \frac{\partial t}{\partial y},\]
e:
\[\frac{dt}{dz} = \sqrt{\frac{a^{2}+s}{b^{2}+s}} \, \frac{1}{D}, \quad \frac{dt}{dy} = \sqrt{\frac{b^{2}+s}{a^{2}+s}} \, \frac{1}{D}.\]

\dots Dunque si trova per integrazione parziale:
\[X = \dots\]
Se si prende la via dell' integrazione come nella espressione di $Z$ il valore dell' integrale sodisfa sempre alla condizione:
\[\frac{\partial X}{\partial z} = \frac{\partial Z}{\partial x},\]
ma pu\`o differire di funzioni di $x$ e di $y$, la funzione sotto segno integrale essendo discontinua anche per $\xi = 0$. Dunque occorre una determinazione ulteriore della via dell' integrazione.

Nella espressione di $\frac{\partial X}{\partial z}$ la funzione sotto segno integrale \`e continua per $s=0$; dunque il pezzo del piano degli $s$, per il cui contorno l'integrale \`e esteso, deve contenere $s = \mathfrak{C}$ e pu\`o contenere o no $s = 0$, ma nessuno altro dei valori sopra notati. Nella espressione di $X$ questo pezzo deve essere determinato in modo che $X$ sia $=0$ per $z=0$; e affinch\`e ci\`o avvenga, dovendo contenere $s = \mathfrak{C}$, deve anche contenere la maggiore radice di $\mathfrak{ts}=0$ (la qu\`ale \`e la maggiore radice di $\mathfrak{t}=0$, se
\[\frac{x^{2}}{a^{2}} + \frac{y^{2}}{b^{2}} - 1 > 0,\]
ed \`e $=0$, se:
\[\frac{x^{2}}{a^{2}} + \frac{y^{2}}{b^{2}} - 1 < 0)\]
ma nessun altra radice di $\mathfrak{ts}=0$. Perch\`e per $\xi=0$ le radici di $F=0$ coincidono colle radici di $\mathfrak{ts}=0$, e se la via dell' integrazione passasse tra due valori di discontinuit\`a che coincidono per $\xi=0$, doverebbe per $\xi=0$ passare per questo valore in modo che l'integrale nella espressione di $X$ diverrebbe infinito ed il valore nonostante il fattore $z$ rimarrebbe finito. ---

Vostro aff.\textsuperscript{mo} Amico Riemann.

\newpage
%==============================================================================
% XVII. Ueber die Flaeche vom kleinsten Inhalt bei gegebener Begrenzung
%==============================================================================
\selectlanguage{german}
\section*{XVII. Ueber die Fl\"ache vom kleinsten Inhalt bei gegebener Begrenzung.\textsuperscript{*})}
\addcontentsline{toc}{section}{XVII. Ueber die Fl\"ache vom kleinsten Inhalt bei gegebener Begrenzung.}

\subsection*{1.}

Eine Fl\"ache l\"asst sich im Sinne der analytischen Geometrie darstellen, indem man die rechtwinkligen Coordinaten $x$, $y$, $z$ eines in ihr beweglichen Punktes als eindeutige Functionen von zwei unabh\"angigen ver\"anderlichen Gr\"ossen $p$ und $q$ angiebt. Nehmen dann $p$ und $q$ bestimmte constante Werthe an, so entspricht dieser einen Combination immer nur ein einziger Punkt der Fl\"ache. Die unabh\"angigen Variabeln $p$ und $q$ k\"onnen in sehr mannigfacher Weise gew\"ahlt werden. F\"ur eine einfach zusammenh\"angende Fl\"ache geschieht dies zweckm\"assig wie folgt. Man l\"asst die Fl\"ache l\"angs der ganzen Begrenzung abnehmen um einen Fl\"achenstreifen, dessen Breite ueberall unendlich klein in derselben Ordnung ist. Durch Wiederholung dieses Verfahrens wird die Fl\"ache fortw\"ahrend verkleinert, bis sie in einen Punkt uebergeht. Die hierbei der Reihe nach auftretenden Begrenzungscurven sind in sich zur\"ucklaufende, von einander getrennte Linien. Man kann sie dadurch unterscheiden, dass man in jeder von ihnen der Gr\"osse $p$ einen besondern constanten Werth beilegt, der um ein Unendlichkleines zu- oder abnimmt, je nachdem man zu der benachbarten umschliessenden oder umschlossenen Curve uebergeht. Die Function $p$ hat dann einen constanten Maximalwerth in der Begrenzung der Fl\"ache und einen Minimalwerth in dem einen Punkte im Innern, in welchen die

\bigskip
\noindent\small\textsuperscript{*}) Dieser Abhandlung liegt ein Manuscript Riemann's zu Grunde, welches nach der eigenen Aeusserung des Verfassers in den Jahren 1860 und 1861 entstanden ist. Dieses Manuscript, welches in gedr\"angter K\"urze nur die Formeln und keinen Text enth\"alt, wurde mir von Riemann im April 1866 zur Bearbeitung anvertraut. Es ist daraus die Abhandlung hervorgegangen, welche ich am 6.\ Januar 1867 der K\"oniglichen Gesellschaft der Wissenschaften zu G\"ottingen eingereicht habe, und welche im 13.\ Band der Abhandlungen dieser Gesellschaft abgedruckt ist. Diese Abhandlung kommt hier in sorgf\"altiger Ueberarbeitung zum zweiten Male zum Abdruck. K.\ Hattendorff.

\newpage
\noindent allm\"ahlich abnehmende Fl\"ache zuletzt zusammenschrumpft. Den Uebergang von einer Begrenzung der abnehmenden Fl\"ache zur n\"achsten kann man dadurch hergestellt denken, dass man jeden Punkt der Curve $(p)$ in einen bestimmten unendlich nahen Punkt der Curve $(p+dp)$ uebergehen l\"asst. Die Wege der einzelnen Punkte bilden dann ein zweites System von Curven, die von dem Punkte des Minimalwerthes von $p$ strahlenf\"ormig nach der Begrenzung der Fl\"ache verlaufen. In jeder dieser Curven legt man $q$ einen besondern constanten Werth bei, der in einer beliebig gew\"ahlten Anfangscurve am kleinsten ist und von da beim Uebergange von einer Curve des zweiten Systems zur andern stetig w\"achst, wenn man zum Zweck des Ueberganges irgend eine Curve $(p)$ in bestimmter Richtung durchl\"auft. Beim Uebergange von der letzten Curve $(q)$ zur Anfangscurve \"andert sich $q$ sprungweise um eine endliche Constante.

Um eine mehrfach zusammenh\"angende Fl\"ache ebenso zu behandeln, kann man sie zuvor durch Querschnitte in eine einfach zusammenh\"angende zerlegen.

Irgend ein Punkt der Fl\"ache l\"asst sich hiernach als Durchschnitt einer bestimmten Curve des Systems $(p)$ mit einer bestimmten Curve des Systems $(q)$ auffassen. Die in dem Punkte $(p,q)$ errichtete Normale verl\"auft von der Fl\"ache aus in zwei entgegengesetzten Richtungen, der positiven und der negativen. Zu ihrer Unterscheidung hat man ueber die gegenseitige Lage der wachsenden positiven Normale, der wachsenden $p$ und der wachsenden $q$ eine Bestimmung zu treffen. Ist nichts anderes festgesetzt, so m\"oge, von der positiven $x$-Axe aus gesehen, die positive $y$-Axe auf dem k\"urzesten Wege in die positive $z$-Axe uebergef\"uhrt werden durch eine Drehung von rechts nach links. Und die Richtung der wachsenden positiven Normale liege zu den Richtungen der wachsenden $p$ und der wachsenden $q$, wie die positive $x$-Axe zur positiven $y$-Axe und zur positiven $z$-Axe. Die Seite der Fl\"ache, auf welcher die positive Normale liegt, soll die positive Seite der Fl\"ache genannt werden.

\subsection*{2.}

Ueber das Gebiet der Fl\"ache sei ein Integral zu erstrecken, dessen Element gleich ist dem Element $dp\,dq$ multiplicirt in eine Functionaldeterminante, also
\[\iint f \frac{\partial(g,h)}{\partial(p,q)}\, dp\, dq,\]
wof\"ur zur Abk\"urzung geschrieben werden soll
\[\iint f \,(dg\, dh).\]

\newpage
Denkt man sich $f$ und $g$ als unabh\"angige Variable eingef\"uhrt, so geht das Integral ueber in $\iint df\, dg$, und es l\"asst sich die Integration nach $f$ oder nach $g$ ausf\"uhren. Die wirkliche Einsetzung von $f$ und $g$ als unabh\"angigen Variabeln verursacht aber Schwierigkeiten oder wenigstens weitl\"aufige Unterscheidungen, wenn dieselbe Werthecombination von $f$ und $g$ in mehreren Punkten der Fl\"ache oder in einer Linie vorhanden ist. Sie ist ganz unm\"oglich, wenn $f$ und $g$ complex sind.

Es ist daher zweckm\"assig, zur Ausf\"uhrung der Integration nach $f$ oder $g$ das Verfahren von Jacobi (Crelle's Journal Bd.~27 p.~208) anzuwenden, bei welchem $p$ und $q$ als unabh\"angige Variable beibehalten werden. Um in Beziehung auf $f$ zu integriren, hat man die Functionaldeterminante in die Form zu bringen
\[\frac{\partial(f,h)}{\partial(p,q)} = \frac{\partial}{\partial q} \left(\frac{\partial f}{\partial p}\, h\right) - \frac{\partial}{\partial p} \left(\frac{\partial f}{\partial q}\, h\right)\]
und erh\"alt zun\"achst
\[\int\!\int \frac{\partial(f,h)}{\partial(p,q)}\, dp\, dq = 0,\]
weil die Integration durch eine in sich zur\"ucklaufende Linie erstreckt wird. Dagegen ist
\[\int\!\int \frac{\partial(g,h)}{\partial(p,q)}\, dp\, dq = \int\!\int \frac{\partial}{\partial q} \left(\frac{\partial g}{\partial p}\, h\right) dp\, dq = \int \frac{\partial g}{\partial p}\, h\, dp\]
in der Richtung der wachsenden $p$ zu nehmen, d.\ h.\ von dem Minimalpunkte im Innern durch eine Curve $(q)$ bis zur Begrenzung. Man erh\"alt $\int h\, dg$ und zwar den Werth, den dieser Ausdruck in der Begrenzung annimmt, da an der untern Grenze des Integrals $h=0$ ist. Folglich wird
\[\int\!\int f \,(dg\, dh) = \int\!\int g \,(df\, dh) = \int g\, dh\]
und das einfache Integral rechts ist in der Richtung der wachsenden $q$ durch die Begrenzung erstreckt. Andererseits hat man nach der eingef\"uhrten Bezeichnung $(df\, dg) = -(dg\, df)$, und daher
\[\int\!\int f \,(df\, dg) = -\int\!\int f \,(dg\, df) = -\int f\, dg,\]
wobei das einfache Integral rechts ebenfalls in der Richtung der wachsenden $q$ durch die Begrenzung der Fl\"ache zu nehmen ist.

\subsection*{3.}

Die Fl\"ache, deren Punkte durch die Curvensysteme $(p)$, $(q)$ festgelegt sind, soll in der folgenden Weise auf einer Kugel vom Radius $1$ abgebildet werden. Im Punkte $(p,q)$ der Fl\"ache, dessen rechtwinklige Coordinaten $x$, $y$, $z$ sind, ziehe man die positive Normale und lege zu ihr eine Parallele durch den Mittelpunkt der Kugel. Der Endpunkt dieser Parallelen auf der Kugeloberfl\"ache ist die Abbildung des Punktes $(x,y,z)$. Durchl\"auft der Punkt $(x,y,z)$ auf der stetig gekr\"ummten Fl\"ache eine zusammenh\"angende Linie, so wird auch die Abbildung derselben auf der Kugel eine zusammenh\"angende Linie sein. Auf dieselbe Weise erh\"alt man als Abbildung eines Fl\"achenst\"ucks ein Fl\"achenst\"uck, als Abbildung der ganzen Fl\"ache eine Fl\"ache, welche die Kugel oder einen Theil derselben einfach oder mehrfach bedeckt.

Der Punkt auf der Kugel, welcher die Richtung der positiven $x$-Axe angiebt, werde zum Pol gew\"ahlt und der Anfangsmeridian durch den Punkt gelegt, welcher der positiven $y$-Axe entspricht. Die Abbildung des Punktes $(x,y,z)$ wird dann auf der Kugel festgelegt durch ihre Poldistanz $r$ und den Winkel $\varphi$, welchen ihr Meridian mit dem Anfangsmeridian einschliesst. F\"ur das Vorzeichen von $\varphi$ gilt die Bestimmung, dass der der positiven $z$-Axe entsprechende Punkt die Coordinaten $r=\frac{\pi}{2}$, $\varphi = +\frac{\pi}{2}$ haben soll.

\subsection*{4.}

Hiernach erh\"alt man als Differential-Gleichung der Fl\"ache
\begin{equation}\tag{1}
\cos r\, dx + \sin r \cos\varphi\, dy + \sin r \sin\varphi\, dz = 0.
\end{equation}
Sind $y$ und $z$ die unabh\"angigen Variabeln, so ergeben sich f\"ur $r$ und $\varphi$ die Gleichungen
\[\cos r = \pm \frac{1}{\sqrt{1 + \left(\frac{\partial x}{\partial y}\right)^{2} + \left(\frac{\partial x}{\partial z}\right)^{2}}},\]
\[\sin r \cos\varphi = \frac{\frac{\partial x}{\partial y}}{\sqrt{1 + \left(\frac{\partial x}{\partial y}\right)^{2} + \left(\frac{\partial x}{\partial z}\right)^{2}}},\]
\[\sin r \sin\varphi = \frac{\frac{\partial x}{\partial z}}{\sqrt{1 + \left(\frac{\partial x}{\partial y}\right)^{2} + \left(\frac{\partial x}{\partial z}\right)^{2}}},\]
in welchen gleichzeitig entweder die oberen oder die unteren Vorzeichen gelten.

\newpage
Ein Parallelogramm auf der positiven Seite der Fl\"ache, begrenzt von den Curven $(p)$ und $(p+dp)$, $(q)$ und $(q+dq)$, projicirt sich auf der $yz$-Ebene in einem Fl\"achenelemente, dessen Inhalt gleich dem absoluten Werthe von $(dy\, dz)$ ist. Das Vorzeichen dieser Functionaldeterminante ist verschieden, je nachdem die im Punkte $(p,q)$ errichtete positive Normale mit der positiven $x$-Axe einen spitzen oder stumpfen Winkel einschliesst. In dem ersten Falle liegen n\"amlich die Projectionen von $dp$ und $dq$ in der $yz$-Ebene ebenso zu einander wie die positive $y$-Axe zur positiven $z$-Axe, im zweiten Falle umgekehrt. Daher ist die Functionaldeterminante im ersten Falle positiv, im zweiten negativ. Und der Ausdruck
\[\sqrt{(dy\, dz)^{2}}\]
ist immer positiv. Er giebt den Inhalt des unendlich kleinen Parallelogramms auf der Fl\"ache. Um also den Inhalt der Fl\"ache selbst zu erhalten, hat man das Doppelintegral
\[\iint \sqrt{(dy\, dz)^{2}}\]
\"uber die ganze Fl\"ache zu erstrecken.

Soll dieser Inhalt ein Minimum sein, so ist die erste Variation des Doppelintegrals $=0$ zu setzen. Man erh\"alt
\[\delta \iint \sqrt{(dy\, dz)^{2}} = \iint \frac{(dy\, dz)\, (d\delta y\, dz + dy\, d\delta z)}{\sqrt{(dy\, dz)^{2}}} = 0,\]
und es gilt das obere oder das untere Zeichen vor der Wurzel, je nachdem $(dy\, dz)$ positiv oder negativ ist. Die linke Seite l\"asst sich schreiben
\[\delta \iint (\sin r \sin\varphi\, dy - \sin r \cos\varphi\, dz)\, (dy\, dz)\]
\[= \iint \delta(\sin r \sin\varphi)\, (dy\, dz) - \iint \delta(\sin r \cos\varphi)\, (dy\, dz)\]
\[+ \iint \sin r \sin\varphi\, (d\delta y\, dz) - \iint \sin r \cos\varphi\, (dy\, d\delta z).\]
Die beiden ersten Integrale reduciren sich auf einfache Integrale, die in der Richtung der wachsenden $q$ durch die Begrenzung der Fl\"ache zu nehmen sind, n\"amlich
\[\int \delta x \, (-\sin r \cos\varphi\, dz + \sin r \sin\varphi\, dy).\]
Der Werth ist $=0$, da in der Begrenzung $\delta x = 0$ ist. Die Bedingung des Minimum lautet also
\[\iint \frac{\partial(\sin r \sin\varphi)}{\partial y} \delta x\, (dy\, dz) + \iint \frac{\partial(\sin r \cos\varphi)}{\partial z} \delta x\, (dy\, dz) = 0.\]
Sie wird erf\"ullt, wenn
\begin{equation}\tag{2}
-\sin r \sin\varphi\, dy + \sin r \cos\varphi\, dz = d\psi
\end{equation}
ein vollst\"andiges Differential ist.

\subsection*{5.}

Die Coordinaten $r$ und $\varphi$ auf der Kugel lassen sich ersetzen durch eine complexe Gr\"osse $\eta = \tg\frac{r}{2} e^{\varphi i}$, deren geometrische Bedeutung leicht zu erkennen ist. Legt man n\"amlich an die Kugel im Pol eine Tangentialebene, deren positive Seite von der Kugel abgekehrt ist, und zieht vom Gegenpol eine Gerade durch den Punkt $(r,\varphi)$, so trifft diese die Tangentialebene in einem Punkte, der die complexe Gr\"osse $2\eta$ repr\"asentirt. Dem Pol entspricht $\eta = 0$, dem Gegenpol $\eta = \infty$. F\"ur die Punkte, welche die Richtungen der positiven $y$- und der positiven $z$-Axe angeben, ist $i\eta = +1$ und resp.\ $=+i$.

F\"uhrt man noch die complexen Gr\"ossen
\[\eta' = \tg\frac{r}{2} e^{-\varphi i}, \quad s = y + zi, \quad s' = y - zi\]
ein, so gehen die Gleichungen (1) und (2) ueber in folgende:
\begin{equation}\tag{1$^{*}$}
(1 - \eta\eta')\, dx + \eta\, ds' + \eta'\, ds = 0,
\end{equation}
\begin{equation}\tag{2$^{*}$}
(1 + \eta\eta')\, d\psi i - \eta'\, ds + \eta\, ds' = 0.
\end{equation}
Diese lassen sich durch Addition und Subtraction verbinden. Dabei werde
\[x + \psi i = 2X, \quad x - \psi i = 2X'\]
gesetzt, so dass umgekehrt $x = X + X'$ ist. Das Problem findet dann seinen analytischen Ausdruck in den beiden Gleichungen
\begin{equation}\tag{3}
ds - \eta\, dX + \eta'\, dX' = 0,
\end{equation}
\begin{equation}\tag{4}
ds' - \eta'\, dX + \eta\, dX' = 0.
\end{equation}
Betrachtet man $X$ und $X'$ als unabh\"angige Variable und stellt die Bedingungen daf\"ur auf, das $ds$ und $ds'$ vollst\"andige Differentiale sind, so findet sich
\[\frac{\partial \eta}{\partial X'} = 0, \quad \frac{\partial \eta'}{\partial X} = 0,\]
d.\ h.\ es ist $\eta$ nur von $X$, $\eta'$ nur von $X'$ abh\"angig, und deshalb umgekehrt $X$ eine Function nur von $\eta$, $X'$ eine Function nur von $\eta'$.

\newpage
Hiernach ist die Aufgabe darauf zur\"uckgef\"uhrt, $\eta$ als Function der complexen Variabeln $X$ oder umgekehrt $X$ als Function der complexen Variabeln $\eta$ so zu bestimmen, dass zugleich den Grenzbedingungen Gen\"uge geleistet werde. Kennt man $\eta$ als Function von $X$, so ergiebt sich daraus $\eta'$, indem man in dem Ausdrucke von $\eta$ jede complexe Zahl in die conjugirte verwandelt. Alsdann hat man nur noch die Gleichungen (3) und (4) zu integriren, um die Ausdr\"ucke f\"ur $s$ und $s'$ zu erlangen. Aus diesen erh\"alt man endlich durch Elimination von $i$ eine Gleichung zwischen $x$, $y$, $z$, die Gleichung der Minimalfl\"ache.

\subsection*{6.}

Sind die Gleichungen (3) und (4) integrirt, so l\"asst sich auch der Inhalt der Minimalfl\"ache selbst leicht angeben, n\"amlich
\[S = 2 \iint \sqrt{(dy\, dz)^{2}}.\]
Die Functionaldeterminante $(dy\, dz)$ formt sich in folgender Weise um
\[(dy\, dz) = \left(\frac{\partial y}{\partial X} \frac{\partial z}{\partial X'} - \frac{\partial y}{\partial X'} \frac{\partial z}{\partial X}\right) (dX\, dX') = -(dX\, dX').\]
Danach erh\"alt man
\[S = 2 \iint (dX\, dX').\]
Zur weiteren Umformung dieses Ausdruckes kann man $y$ aus $Y$ und $Y'$, $z$ aus $Z$ und $Z'$ ebenso zusammensetzen wie $x$ aus $X$ und $X'$, so dass die Gleichungen gelten
\[x = X + X', \quad x' = X - X',\]
\[y = Y + Y', \quad y' = Y - Y',\]
\[z = Z + Z', \quad z' = Z - Z'.\quad \text{q.\ e.\ d.}\]

\newpage
Alsdann erh\"alt man schliesslich
\begin{equation}\tag{5}
\begin{aligned}
S &= -i \iint \left[(dX\, dX') + (dY\, dY') + (dZ\, dZ')\right] \\
&= \tfrac{1}{2} \iint \left[(dx\, dx') + (dy\, dy') + (dz\, dz')\right].
\end{aligned}
\end{equation}

\subsection*{7.}

Die Minimalfl\"ache und ihre Abbildungen auf der Kugel wie in den Ebenen, deren Punkte resp.\ die complexen Gr\"ossen $\eta$, $X$, $Y$, $Z$ repr\"asentiren, sind einander in den kleinsten Theilen \"ahnlich. Man erkennt dies sofort, wenn man das Quadrat des Linearelementes in diesen Fl\"achen ausdr\"uckt. Dasselbe ist

auf der Kugel $\sin r^{2}\, d\log\eta\, d\log\eta'$,

in der Ebene der $\eta$ $\frac{4}{(1+\eta\eta')^{2}}\, d\eta\, d\eta'$,

in der Ebene der $X$ $\frac{dX}{d\eta} \frac{dX'}{d\eta'}\, d\eta\, d\eta'$,

in der Ebene der $Y$ $\frac{dY}{d\eta} \frac{dY'}{d\eta'}\, d\eta\, d\eta'$,

in der Ebene der $Z$ $\frac{dZ}{d\eta} \frac{dZ'}{d\eta'}\, d\eta\, d\eta'$,

in der Minimalfl\"ache selbst
\[\begin{aligned}
dx^{2} + dy^{2} + dz^{2} &= (dX+dX')^{2} + (dY+dY')^{2} + (dZ+dZ')^{2} \\
&= 2(dX\, dX' + dY\, dY' + dZ\, dZ') \\
&= 2\left(\frac{dX}{d\eta} \frac{dX'}{d\eta'} + \frac{dY}{d\eta} \frac{dY'}{d\eta'} + \frac{dZ}{d\eta} \frac{dZ'}{d\eta'}\right) d\eta\, d\eta'.
\end{aligned}\]
Es ist n\"amlich nach den Gleichungen (3) und (4), wenn man darin $\eta$ und $\eta'$ als unabh\"angige Variable ansieht:
\[\frac{dX}{d\eta} = \frac{ds}{d\eta} \frac{1}{2\eta'}, \quad \frac{dX'}{d\eta'} = \frac{ds'}{d\eta'} \frac{1}{2\eta},\]
\[\frac{dX}{d\eta'} = \frac{ds}{d\eta'} \frac{1}{2\eta}, \quad \frac{dX'}{d\eta} = \frac{ds'}{d\eta} \frac{1}{2\eta'},\]
und deshalb
\[dX^{2} + dY^{2} + dZ^{2} = 0,\]
\[dX'^{2} + dY'^{2} + dZ'^{2} = 0.\]
Das Verh\"altniss von irgend zwei der obigen Quadrate von Linearelementen ist unabh\"angig von $d\eta$ und $d\eta'$, d.\ h.\ von der Richtung des Elementes, und darin beruht die in den kleinsten Theilen \"ahnliche Abbildung. Da die Linearvergr\"osserung bei der Abbildung in irgend einem Punkte nach allen Richtungen dieselbe ist, so erh\"alt man die Fl\"achenvergr\"osserung gleich dem Quadrat der Linearvergr\"osserung. Das Quadrat des Linearelementes in der Minimalfl\"ache ist aber gleich

\newpage
\noindent der doppelten Summe der Quadrate der entsprechenden Linearelemente in den Ebenen der $X$, der $Y$ und der $Z$. Daher ist auch das Fl\"achenelement in der Minimalfl\"ache gleich der doppelten Summe der entsprechenden Fl\"achenelemente in jenen Ebenen. Dasselbe gilt von der ganzen Fl\"ache und ihren Abbildungen in den Ebenen der $X$, $Y$, $Z$.

\subsection*{8.}

Eine wichtige Folgerung l\"asst sich noch aus dem Satze von der Aehnlichkeit in den kleinsten Theilen ziehen, wenn man eine neue complexe Variable $\eta_{1}$ dadurch einf\"uhrt, dass man auf der Kugel den Pol in einen beliebigen Punkt $(\eta = a)$ verlegt und den Anfangsmeridian beliebig w\"ahlt. Hat dann $\eta_{1}$ f\"ur das neue Coordinatensystem dieselbe Bedeutung wie $\eta$ f\"ur das alte, so kann man jetzt ein unendlich kleines Dreieck auf der Kugel sowohl in der Ebene der $\eta$ als in der der $\eta_{1}$ abbilden. Die beiden Bilder sind dann auch Abbildungen von einander und in den kleinsten Theilen \"ahnlich. F\"ur den Fall der directen Aehnlichkeit ergiebt sich ohne Weiteres, dass $\frac{d\eta_{1}}{d\eta}$ unabh\"angig ist von der Richtung der Verschiebung von $\eta$, d.\ h.\ dass $\eta_{1}$ eine Function der complexen Variabeln $\eta$ ist. Den Fall der inversen (symmetrischen) Aehnlichkeit kann man auf den vorigen zur\"uckf\"uhren, indem man statt $\eta_{1}$ die conjugirte complexe Gr\"osse nimmt. Um nun $\eta_{1}$ als Function von $\eta$ auszudr\"ucken, hat man zu beachten, dass $\eta_{1} = 0$ ist in dem einen Punkte der Kugel, f\"ur welchen $\eta = a$, und $\eta_{1} = \infty$ in dem diametral gegen\"uberliegenden Punkte, d.\ h.\ f\"ur $\eta = -\frac{1}{a'}$. Danach ergiebt sich $\eta_{1} = c\frac{\eta - a}{1 + a'\eta}$. Zur Bestimmung der Constanten $c$ dient die Bemerkung, dass, wenn $\eta_{1} = \beta$ ist f\"ur $\eta = 0$, daraus $\eta_{1} = -\frac{1}{\beta'}$ gefunden wird f\"ur $\eta = \infty$. Es ist also $\beta = -ca'$ und $-\frac{1}{\beta'} = \frac{c}{a}$, d.\ h.\ $\beta\beta' = \frac{c}{a}a'$. Hieraus ergiebt sich $cc' = 1$ und daher $c = e^{\vartheta i}$ f\"ur ein reelles $\vartheta$. Die Gr\"ossen $\alpha$ und $\vartheta$ k\"onnen beliebige Werthe erhalten: $\alpha$ h\"angt von der Lage des neuen Pols, $\vartheta$ von der Lage des neuen Anfangsmeridians ab. Diesem neuen Coordinatensystem auf der Kugel entsprechen die Richtungen der Axen eines neuen rechtwinkligen Systems. Es m\"ogen in dem neuen System $X_{1}$, $s_{1}$, $s'_{1}$ dasselbe bezeichnen wie $X$, $s$, $s'$ in dem alten. Dann erlangt man die Transformationsformeln

\newpage
\begin{equation}\tag{6}
\begin{aligned}
(1 + \alpha\alpha') X_{1} &= (1 - \alpha\alpha') X + \alpha' s + \alpha s', \\
(1 + \alpha\alpha') s'_{1} e^{\vartheta i} &= -2\alpha X - \alpha^{2} s + s.
\end{aligned}
\end{equation}

\subsection*{9.}

Aus den Transformationsformeln (6) berechnen wir
\[\left(\frac{d\eta_{1}}{d\eta}\right)^{2} \frac{dX_{1}}{d\eta_{1}} = \frac{dX}{d\eta}\]
oder
\[\left(\frac{d\eta_{1}}{d\eta}\right)^{2} dX_{1} = \left(\frac{d\log\eta_{1}}{d\log\eta}\right)^{2} dX.\]
Hiernach empfiehlt es sich, eine neue complexe Gr\"osse $u$ einzuf\"uhren, welche durch die Gleichung definirt wird
\begin{equation}\tag{7}
du = \left(\frac{d\log\eta_{1}}{d\log\eta}\right)^{2} dX,
\end{equation}
und die von der Lage des Coordinatensystems $(x,y,z)$ unabh\"angig ist. Gelingt es dann, $u$ als Function von $\eta$ zu bestimmen, so erh\"alt man
\begin{equation}\tag{8}
X = \int du \left(\frac{d\log\eta_{1}}{d\log\eta}\right)^{-2}.
\end{equation}
$X$ ist der Abstand des zu $\eta$ geh\"origen Punktes der Minimalfl\"ache von einer Ebene, die durch den Anfangspunkt der Coordinaten rechtwinklig zur Richtung $\eta = 0$ gelegt ist. Man erh\"alt den Abstand desselben Punktes der Minimalfl\"ache von einer durch den Anfangspunkt der Coordinaten gelegten Ebene, die rechtwinklig auf der Richtung $\eta = \alpha$ steht, indem man in (8) $-e^{\vartheta i}\frac{1}{\eta}$ statt $\eta$ setzt. Speciell also f\"ur $\alpha = 1$ und $\alpha = i$
\begin{equation}\tag{9}
x = \int du, \quad y = \int \frac{du}{\eta^{2}}, \quad z = \int \frac{du}{\eta^{2}}.
\end{equation}

\newpage
\subsection*{10.}

Die Gr\"osse $u$ ist als Function von $\eta$ zu bestimmen, d.\ h.\ als einwerthige Function des Ortes in derjenigen Fl\"ache, welche, ueber die $\eta$-Ebene ausgebreitet, die Minimalfl\"ache in den kleinsten Theilen \"ahnlich abbildet. Daher kommt es vor allen Dingen auf die Unstetigkeiten und Verzweigungen in dieser Abbildung an. Bei der Untersuchung derselben hat man Punkte im Innern der Fl\"ache von Begrenzungspunkten zu unterscheiden.

Handelt es sich um einen Punkt im Innern der Minimalfl\"ache, so lege man in ihn den Anfangspunkt des Coordinatensystems $(x,y,z)$, die Axe der positiven $x$ in die positive Normale, folglich die $yz$-Ebene tangential. Dann fehlen in der Entwicklung von $x$ das freie Glied und die in $y$ und $z$ multiplicirten Glieder. Durch geeignet gew\"ahlte Richtung der $y$-Axe und der $z$-Axe kann man auch das in $yz$ multiplicirte Glied verschwinden lassen. Die partielle Differentialgleichung der Minimalfl\"ache reducirt sich unter dieser Voraussetzung f\"ur unendlich kleine Werthe von $y$ und $z$ auf $\frac{\partial^{2}x}{\partial y^{2}} + \frac{\partial^{2}x}{\partial z^{2}} = 0$. Das Kr\"ummungsmass ist also negativ, die Haupt-Kr\"ummungsradien sind einander entgegengesetzt gleich. Die Tangentialebene theilt die Fl\"ache in vier Quadranten, wenn die Kr\"ummungshalbmesser nicht $\infty$ sind. Diese Quadranten liegen abwechselnd ueber und unter der Tangentialebene. Beginnt die Entwicklung von $x$ erst mit den Gliedern $n$ter Ordnung ($n>2$), so sind die Kr\"ummungsradien $\infty$, und die Tangentialebene theilt die Fl\"ache in $2n$ Sectoren, die abwechselnd ueber und unter jener Ebene liegen und von den Kr\"ummungslinien halbirt werden.\textsuperscript{(*)})

Will man nun $X$ als Function der complexen Variabeln $Y$ ansehen, so ergiebt sich in dem Falle der vier Sectoren
\[\log X = 2\log Y + \text{funct.\ cont.},\]
in dem Falle der $2n$ Sectoren
\[\log X = n\log Y + \text{f.\ c.}\]
Und da nach (8) und (9) $\frac{du}{d\log\eta} = dX$ ist, so beginnt die Entwicklung von $\eta$ im ersten Falle mit der ersten, im zweiten mit der $(n-1)$ten Potenz von $Y$. Umgekehrt wird also, wenn $Y$ als Function von $\eta$ angesehen werden soll, die Entwicklung im ersten Falle nach ganzen Potenzen von $\eta$, im zweiten nach ganzen Potenzen von $\eta^{\frac{1}{n-1}}$ fortschreiten. D.\ h.\ die Abbildung auf der $\eta$-Ebene hat an der betreffenden Stelle keinen oder einen $(n-2)$fachen Verzweigungspunkt, je nachdem der erste oder der zweite Fall eintritt.

\newpage
Was $u$ betrifft, so ergiebt sich $\frac{du}{d\log Y} = \frac{du}{d\log\eta} \frac{d\log\eta}{d\log Y}$, also mit H\"ulfe der Gleichung (9)
\[\frac{du}{d\log Y} = 2\, \frac{dX}{d\log\eta}\, \frac{d\log\eta}{d\log Y}.\]
Demnach ist in einem $(n-2)$fachen Verzweigungspunkte der Abbildung auf der $\eta$-Ebene
\[\log \frac{du}{d\log Y} = 2\log Y + \text{f.\ c.}\]
oder
\[u = \text{const.} + Y^{2}.\]

\subsection*{11.}

Die weitere Untersuchung soll zun\"achst auf den Fall beschr\"ankt werden, dass die gegebene Begrenzung aus geraden Linien besteht. Dann l\"asst sich die Abbildung der Begrenzung auf der $\eta$-Ebene wirklich herstellen. Die in irgend welchen Punkten einer geraden Begrenzungslinie errichteten Normalen liegen in parallelen Ebenen, und daher ist die Abbildung auf der Kugel ein Bogen eines gr\"ossten Kreises.

Um einen Punkt im Innern einer geraden Begrenzungslinie zu untersuchen, legt man wie vorher in ihn den Anfangspunkt der Coordinaten, die positive $x$-Axe in die positive Normale. Dann f\"allt die ganze Begrenzungslinie in die $yz$-Ebene. Der reelle Theil von $X$ ist demnach in der ganzen Begrenzungslinie $=0$. Geht man also durch das Innere der Minimalfl\"ache um den Anfangspunkt der Coordinaten herum von einem vorangehenden bis zu einem nachfolgenden Begrenzungspunkte, so muss dabei der Arcus von $X$ sich \"andern um $n\pi$, ein ganzes Vielfaches von $\pi$. Der Arcus von $Y$ \"andert sich gleichzeitig um $\pi$. Man hat also, wie vorher
\[\log X = n\log Y + \text{f.\ c.}\]
\[\log \eta = (n-1)\log Y + \text{f.\ c.}\]
Dem betrachteten Begrenzungspunkte entspricht ein $(n-2)$facher Verzweigungspunkt in der Abbildung auf der $\eta$-Ebene. In dieser Abbildung macht das auf den Punkt folgende Begrenzungsst\"uck mit dem ihm vorhergehenden den Winkel $(n-1)\pi$.

\subsection*{12.}

Bei dem Uebergange von einer Begrenzungslinie zur folgenden hat man zwei F\"alle zu unterscheiden. Entweder treffen sie zusammen in einem im Endlichen liegenden Schnittpunkte, oder sie erstrecken sich ins Unendliche.

Im ersten Falle sei $\alpha\pi$ der im Innern der Minimalfl\"ache liegende Winkel der beiden Begrenzungslinien. Legt man den Anfangspunkt der Coordinaten in den zu untersuchenden Eckpunkt, die positive $x$-Axe in die positive Normale, so ist in beiden Begrenzungslinien der reelle Theil von $X = 0$. Beim Uebergange von der ersten Begrenzungslinie zur folgenden \"andert sich also der Arcus von $X$ um $m\pi$, ein ganzes Vielfaches von $\pi$, der Arcus von $Y$ um $\alpha\pi$. Man hat daher
\[\log X = m\log Y + \text{f.\ c.}\]
\[\left(1 - \frac{\alpha}{m}\right)\log\eta = \log Y + \text{f.\ c.}\]

Erstreckt sich die Fl\"ache zwischen zwei auf einander folgenden Begrenzungsgeraden ins Unendliche, so lege man die positive $x$-Axe in ihre k\"urzeste Verbindungslinie, parallel der positiven Normalen im Unendlichen. Die L\"ange der k\"urzesten Verbindungslinie sei $A$, und $\alpha\pi$ der Winkel, welchen die Projection der Minimalfl\"ache in der $yz$-Ebene ausf\"ullt. Dann bleiben die reellen Theile von $X$ und $i\, d\log\eta$ im Unendlichen endlich und stetig und nehmen in den begrenzenden Geraden constante Werthe an. Hieraus ergiebt sich (f\"ur $y=\infty$, $z=\infty$)
\[X = \frac{A}{2\alpha\pi}\, d\log\eta + \text{f.\ c.}\]
Legt man die $x_{1}$-Axe eines Coordinatensystems in eine begrenzende Gerade, die $x_{2}$-Axe eines andern Systems in die zweite begrenzende Gerade u.\ s.\ f., so ist in der ersten Linie $\log\eta_{1}$, in der zweiten $\log\eta_{2}$ u.\ s.\ f.\ rein imagin\"ar, da die Normale zu der betreffenden Axe der $x_{1}$, der $x_{2}$ u.\ s.\ f.\ senkrecht steht. Es ist also $i\frac{du}{d\log\eta_{1}}$ in der ersten Begrenzungslinie reell, $i\frac{du}{d\log\eta_{2}}$ in der zweiten u.\ s.\ f.\ Da aber auch f\"ur ein beliebiges Coordinatensystem $(x,y,z)$ immer
\[i\, du = i\left(\frac{d\log\eta_{1}}{d\log\eta}\right)^{2} \frac{du}{d\log\eta_{1}}\]
ist, so findet sich, dass in jeder geraden Begrenzungslinie
\[i\sqrt{\frac{du}{d\log\eta}}\]
entweder reelle oder rein imagin\"are Werthe besitzt.

\newpage
\subsection*{13.}

Die Minimalfl\"ache ist bestimmt, sobald man eine der Gr\"ossen $u$, $\eta$, $X$, $Y$, $Z$ durch eine der uebrigen ausgedr\"uckt hat. Dies gelingt in vielen F\"allen. Besondere Beachtung verdienen darunter diejenigen, in welchen $\frac{du}{d\log\eta}$ eine algebraische Function von $\eta$ ist. Dazu ist n\"othig und hinreichend, dass die Abbildung auf der Kugel und ihre symmetrischen und congruenten Fortsetzungen eine geschlossene Fl\"ache bilden, welche die ganze Kugel einfach oder mehrfach bedeckt.

Im Allgemeinen aber wird es schwierig sein, direct eine der Gr\"ossen $u$, $\eta$, $X$, $Y$, $Z$ durch eine der uebrigen auszudr\"ucken. Statt dessen kann man aber auch jede von ihnen als Function einer neuen zweckm\"assig gew\"ahlten unabh\"angigen Variabeln bestimmen. Wir f\"uhren eine solche unabh\"angige Variable $t$ ein, dass die Abbildung der Fl\"ache auf der $t$-Ebene die halbe unendliche Ebene einfach bedeckt, und zwar diejenige H\"alfte, f\"ur welche der imagin\"are Theil von $t$ positiv ist. In der That ist es immer m\"oglich, $t$ als Function von $u$ (oder von irgend einer der uebrigen Gr\"ossen $\eta$, $X$, $Y$, $Z$) in der Fl\"ache so zu bestimmen, dass der imagin\"are Theil in der Begrenzung $=0$ ist, und dass sie in einem beliebigen Begrenzungspunkte ($u=b$) unendlich von der ersten Ordnung wird, d.\ h.
\[t = \frac{\text{const.}}{u-b}.\]
Der Arcus des Factors von $\frac{1}{u-b}$ ist durch die Bedingung bestimmt, dass der imagin\"are Theil von $t$ in der Begrenzung $=0$, im Innern der Fl\"ache positiv sein soll. Es bleibt also in dem Ausdrucke von $t$ nur der Modul dieses Factors und eine additive Constante willk\"urlich.

Es sei $t = a_{1}, a_{2}, \dots$ f\"ur die Verzweigungspunkte im Innern der Abbildung auf der $\eta$-Ebene, $t = b_{1}, b_{2}, \dots$ f\"ur die Verzweigungspunkte in der Begrenzung, die nicht Eckpunkte sind, $t = c_{1}, c_{2}, \dots$ f\"ur die Eckpunkte, $t = e_{1}, e_{2}, \dots$ f\"ur die ins Unendliche sich erstreckenden Sectoren. Wir wollen der Einfachheit wegen voraussetzen, dass die s\"ammtlichen Gr\"ossen $a$, $b$, $c$, $e$ im endlichen Gebiete der $t$-Ebene liegen.

Dann hat man

f\"ur $t = a$: $\log\frac{du}{d\log\eta} = (\gamma - 1)\log(t-a) + \text{f.\ c.}$,

f\"ur $t = b$: $\log\frac{du}{d\log\eta} = \left(\frac{\gamma}{2} - 1\right)\log(t-b) + \text{f.\ c.}$,

f\"ur $t = c$: $\log\frac{du}{d\log\eta} = \left(\frac{\gamma}{2} - 1\right)\log(t-c) + \text{f.\ c.}$,

f\"ur $t = e$: $u = \gamma\sqrt{t-e}\, \log(t-e) + \text{f.\ c.}$

\newpage
Man kann die Untersuchung auf den Fall $n=3$, $m=1$ beschr\"anken, d.\ h.\ auf einfache Verzweigungspunkte, und den allgemeinen Fall aus diesem dadurch ableiten, dass man mehrere einfache Verzweigungspunkte zusammenfallen l\"asst.

Um den Ausdruck f\"ur $\frac{du}{dt}$ zu bilden, hat man zu beachten, dass l\"angs der Begrenzung $dt$ reell, $du$ entweder reell oder rein imagin\"ar ist. Demnach ist $i\frac{du}{dt}$ reell, wenn $t$ reell ist. Diese Function kann man ueber die Linie der reellen Werthe von $t$ hin\"uber stetig fortsetzen, indem man die Bestimmung trifft, dass f\"ur conjugirte Werthe $t$ und $t'$ der Variabeln auch die Function conjugirte Werthe haben soll. Alsdann ist $i\frac{du}{dt}$ f\"ur die ganze $t$-Ebene bestimmt und zeigt sich einwerthig.

Es seien $a'_{1}, a'_{2}, \dots$ die conjugirten Werthe zu $a_{1}, a_{2}, \dots$; und das Product $(t-a_{1})(t-a_{2})\cdots$ werde mit $\prod(t-a)$ bezeichnet. Alsdann ist
\begin{equation}\tag{11}
du = \text{const.} \cdot i \sqrt{ \frac{ \prod(t-a)\prod(t-a')\prod(t-b) }{ \prod(t-c)^{2}\prod(t-e)^{2}} } \, dt.
\end{equation}
Die Constanten $a$, $b$, $c$ etc.\ m\"ussen so bestimmt werden, dass f\"ur $t=e$: $u = \sqrt{t-e}\, \log(t-e) + \text{f.\ c.}$ wird. Damit $u$ f\"ur alle Werthe von $t$ ausser $a$, $b$, $c$, $e$ endlich und stetig bleibe, muss f\"ur die Anzahl dieser letztgenannten Werthe eine Relation bestehen. Es muss die Differenz der Anzahl der Eckpunkte und der in der Begrenzung liegenden Verzweigungspunkte um $4$ gr\"osser sein als die doppelte Differenz der Anzahl der innern Verzweigungspunkte und der ins Unendliche verlaufenden Sectoren. Setzt man zur Abk\"urzung
\[\prod(t-a)\prod(t-a')\prod(t-b) = \varphi(t),\]
\[\prod(t-c)\prod(t-e) = \chi(t),\]
\newpage
\noindent d.\ h.
\[\frac{du}{dv} = \text{const.} \cdot i \sqrt{\frac{\varphi(t)}{\chi(t)^{2}}},\]
so ist die ganze Function $\varphi(t)$ vom Grade $\nu - 4$, wenn $\chi(t)$ vom Grade $\nu$ ist. Hier bedeutet $\nu$ die Anzahl der Eckpunkte vermehrt um die doppelte Anzahl der ins Unendliche verlaufenden Sectoren.

\subsection*{14.}

Es ist noch $\eta$ als Function von $t$ auszudr\"ucken. Direct gelangt man dazu nur in den einfachsten F\"allen. Im Allgemeinen ist der folgende Weg einzuschlagen. Es sei $v$ eine noch n\"aher zu bestimmende Function von $t$, die als bekannt vorausgesetzt wird. In den Gleichungen (8), (9), (10) kommt es wesentlich an auf $\frac{dX}{d\log\eta}$, wof\"ur man auch schreiben kann $\frac{dX}{dv} \cdot \frac{dv}{d\log\eta}$.

Der letzte Factor l\"asst sich ansehen als Product der beiden Factoren
\begin{equation}\tag{12}
k_{1} = \sqrt{\frac{dv}{d\eta}}, \quad k_{2} = \eta\sqrt{\frac{dv}{d\eta}},
\end{equation}
die der Differentialgleichung erster Ordnung gen\"ugen
\begin{equation}\tag{13}
k_{2} \frac{dk_{1}}{dv} - k_{1} \frac{dk_{2}}{dv} = 1,
\end{equation}
sowie der Differentialgleichung zweiter Ordnung
\begin{equation}\tag{14}
\frac{1}{k} \frac{d^{2}k}{dv^{2}} = \frac{1}{k_{1}} \frac{d^{2}k_{1}}{dv^{2}} = \frac{1}{k_{2}} \frac{d^{2}k_{2}}{dv^{2}}.
\end{equation}
Gelingt es also, die eine oder die andere Seite dieser letzten Gleichung als Function von $t$ auszudr\"ucken, so l\"asst sich eine homogene lineare Differentialgleichung zweiter Ordnung herstellen, von welcher $k_{1}$ und $k_{2}$ particul\"are Integrale sind. Es sei $k$ das vollst\"andige Integral. Wir ersetzen $\frac{d^{2}k}{dv^{2}}$ durch das ihm gleichbedeutende
\[\frac{dv}{dt} \frac{d^{2}k}{dt^{2}} - \frac{dk}{dt} \frac{d^{2}v}{dt^{2}} \bigg/ \left(\frac{dv}{dt}\right)^{2}\]
und erhalten f\"ur $k$ die Differentialgleichung
\begin{equation}\tag{15}
\frac{dv}{dt} \frac{d^{2}k}{dt^{2}} - \frac{d^{2}v}{dt^{2}} \frac{dk}{dt} = \left(\frac{dv}{dt}\right)^{3} \left(\frac{1}{k_{1}} \frac{d^{2}k_{1}}{dv^{2}}\right) k.
\end{equation}
Von der Gleichung (15) seien zwei von einander unabh\"angige particul\"are Integrale $K_{1}$ und $K_{2}$ gefunden, deren Quotient $K_{2}:K_{1} = H$

\newpage
\noindent ein von Bogen gr\"osster Kreise begrenztes Abbild der positiven $t$-Halbebene auf der Kugelfl\"ache liefert. Dasselbe leistet dann jeder Ausdruck von der Form
\begin{equation}\tag{16}
\eta = e^{\vartheta i} \frac{K_{2} - \alpha K_{1}}{K_{2} - \alpha' K_{1}},
\end{equation}
worin $\vartheta$ reell und $\alpha$, $\alpha'$ conjugirte complexe Gr\"ossen sind.

Die Function $v$ ist so zu w\"ahlen, dass f\"ur endliche Werthe von $t$ die Unstetigkeiten von $\frac{1}{k} \frac{d^{2}k}{dv^{2}}$ nicht ausserhalb der Punkte $a$, $a'$, $b$, $c$, $e$ liegen.

Setzt man
\begin{equation}\tag{17}
\frac{dv}{dt} = \frac{1}{\sqrt{\psi(t)\chi(t)}} = \frac{1}{\Pi(t)},
\end{equation}
so wird die Function $\frac{1}{k} \frac{d^{2}k}{dv^{2}}$ im Endlichen unstetig nur f\"ur die Punkte $a$, $a'$, $b$, $c$, und zwar f\"ur jeden unendlich in erster Ordnung. Man erh\"alt n\"amlich f\"ur $t = c$
\[\frac{1}{k} \frac{d^{2}k}{dv^{2}} = \frac{\text{const.}}{t-c}.\]
Folglich:
\[k = \text{const.} \cdot (t-c)^{\frac{1}{2}}, \quad \eta = \text{const.} \cdot (t-c),\]
und hieraus:
\[\eta - \eta_{c} = \text{const.} \cdot (t-c)^{\frac{\gamma}{2}}.\]
F\"ur $t = c$:
\[k \frac{dv^{2}}{dt^{2}} : \frac{1}{t-c}.\]
Entsprechende Ausdr\"ucke erh\"alt man f\"ur $t = a$, $a'$, $b$, in denen $c$ resp.\ durch $a$, $a'$, $b$, und $\gamma$ durch $2$ zu ersetzen ist.

Eine \"ahnliche Betrachtung lehrt, dass f\"ur $t = e$ die Function $\frac{1}{k} \frac{d^{2}k}{dv^{2}}$ stetig bleibt.

F\"ur $t = \infty$ ergiebt sich
\[\frac{1}{k} \frac{d^{2}k}{dv^{2}} = \frac{\text{const.}}{t^{2}}.\]
\[\frac{1}{k} \frac{d^{2}k}{dv^{2}} = \sum \frac{\frac{1}{2}(\frac{\gamma}{2}-1)}{t-g} + \frac{F(t)}{\psi(t)\chi(t)}.\]
Die Summe bezieht sich auf alle Punkte $g = a$, $a'$, $b$, $c$, und bei $a$, $a'$, $b$ ist $2$ statt $\gamma$ zu setzen. $F(t)$ ist eine ganze Function vom Grade $(2\nu - 6)$, in der die ersten beiden Coefficienten sich folgendermassen bestimmen. Man bringe $dv$ in die Form

\newpage
\noindent $dv = \sqrt{a}\, dv'$

oder k\"urzer $dv^{2} = a\, dv'^{2}$.

Dann ergiebt sich durch Differentiation
\[2\, dv \frac{d^{2}v}{dt^{2}} = \frac{da}{dt}\, dv'^{2} + 2a\, dv' \frac{d^{2}v'}{dt^{2}},\]
folglich
\[\frac{1}{dv} \frac{d^{2}v}{dt^{2}} = \frac{1}{2} \frac{da}{dt}\frac{1}{a} + \frac{dv'}{dv} \frac{d^{2}v'}{dt^{2}}.\]
\[\left(\frac{d^{2}v}{dv^{2}}\right)^{2} \frac{d}{dv'} \left\{\left(\frac{dv}{dv'}\right)^{2} \frac{d^{2}v}{dt^{2}}\right\} = a^{2} \frac{d^{2}v}{dv^{2}}\]
\[= \sum \frac{\frac{1}{2}(\frac{\gamma}{2}-1)}{t-g} + \frac{F(t)}{\psi(t)\chi(t)}.\]
Die Function auf der linken Seite ist endlich f\"ur $t = \infty$. Folglich hat man rechts in der Entwicklung von $\frac{1}{\psi(t)\chi(t)} F(t)$ und von $\frac{dv}{dt}$ die Coefficienten von $t^{2}$ und resp.\ von $t$ einander gleich zu setzen. Die Entwickelung von $\frac{dv}{dt}$ giebt nach einfacher Rechnung

Hiernach bleiben in $F(t)$ noch $2\nu - 7$ unbestimmte Coefficienten. Es ist aber wichtig zu bemerken, dass dieselben reell sein m\"ussen. Denn wir haben in \S.~12 gefunden, dass $du$ reell oder rein imagin\"ar ist in allen geraden Begrenzungslinien der Minimalfl\"ache und folglich auch an jeder Stelle in der Begrenzung der Abbildungen. Verm\"oge der Gleichung (17) gilt dasselbe von $dv$. Daraus l\"asst sich beweisen, dass f\"ur reelle Werthe von $t$ die Function $\frac{1}{k} \frac{d^{2}k}{dv^{2}}$ nothwendigerweise reelle Werthe besitzt.

Um diesen Beweis zu f\"uhren, betrachten wir die Abbildung auf der Kugel vom Radius $1$ und nehmen irgend einen Theil der Begrenzung, also den Bogen eines gewissen gr\"ossten Kreises. Im Pole dieses gr\"ossten Kreises legen wir die Tangential-Ebene an und bezeichnen

\newpage
\noindent sie als die Ebene der $\eta_{1}$. Dann lassen sich die constanten Gr\"ossen $\alpha_{1}$, $\alpha'_{1}$, $\vartheta_{1}$ so bestimmen, dass
\[\eta_{1} = e^{\vartheta_{1}i} \frac{\eta - \alpha_{1}}{1 + \alpha'_{1}\eta}\]
ist, und wir erhalten zwei Functionen $k'_{1} = \sqrt{\frac{dv}{d\eta_{1}}}$ und $k'_{2} = \eta_{1}\sqrt{\frac{dv}{d\eta_{1}}}$, die particul\"are Integrale der Differentialgleichung (15) sind. Folglich haben wir
\[\frac{1}{k_{1}} \frac{d^{2}k_{1}}{dv^{2}} = \frac{1}{k'_{1}} \frac{d^{2}k'_{1}}{dv^{2}}.\]
Der eben betrachtete Theil der Begrenzung bildet sich in der $\eta_{1}$-Ebene ab durch die Gleichung
\[\eta_{1} + \eta'_{1} = 0\]
und wenn man dies in $k'_{1}$ einf\"uhrt, so erkennt man leicht, dass in dem fraglichen Begrenzungstheile $\frac{1}{k'_{1}} \frac{d^{2}k'_{1}}{dv^{2}}$ reell ausf\"allt. Folglich gilt dasselbe von $\frac{1}{k_{1}} \frac{d^{2}k_{1}}{dv^{2}}$, und da diese Betrachtung f\"ur jedes einzelne Begrenzungsst\"uck angestellt werden kann, so ist $\frac{1}{k} \frac{d^{2}k}{dv^{2}}$ reell in der ganzen Begrenzung.

Nun f\"allt aber bei einem reellen oder rein imagin\"aren $dv$ die Function $\frac{1}{k_{1}} \frac{d^{2}k_{1}}{dv^{2}}$ auch dann reell aus, wenn man allgemeiner
\[k'_{1} = \varrho_{1} e^{\psi_{1}i}\]
setzt und den Modul $\varrho_{1}$ constant nimmt. Damit also die Axe der reellen $t$ sich auf der Kugel vom Radius $1$ wirklich in B\"ogen gr\"osster Kreise abbilde, muss f\"ur jeden Begrenzungstheil $\varrho_{1} = 1$ sein. Dies liefert ebenso viele Bedingungsgleichungen, als einzelne Begrenzungslinien gegeben sind.

Bei dieser Untersuchung ist, wie schon im vorigen Paragraphen, vorausgesetzt, dass die Werthe $a$, $b$, $c$, $e$ s\"ammtlich endlich seien. Trifft dies nicht zu, so bedarf die Betrachtung einer geringen Modification.

\medskip
\noindent\textbf{Anmerkung.} Die Aufgabe ist hiermit vollst\"andig formulirt. Im einzelnen Falle kommt es nur darauf an, die Differentialgleichung (15) wirklich aufzustellen und zu integriren. Uebrigens ist es nicht unwichtig, zu bemerken, dass die Anzahl der in der L\"osung auftretenden willk\"urlichen reellen Constanten ebenso gross ist wie die Anzahl der Bedingungsgleichungen, welche verm\"oge der Natur der Aufgabe und verm\"oge der Daten des Problems erf\"ullt sein m\"ussen. Wir bezeichnen die Anzahl der Punkte $a$, $b$, $c$, $e$ resp.\ mit $A$, $B$, $C$, $E$ und beachten, dass $2A + B + 4 = C + 2E = \nu$ ist. In der Differentialgleichung (15) treten $2A + B + 4C + 5E - 10$ willk\"urliche reelle Constanten auf, n\"amlich: die

\newpage
\noindent Winkel $\gamma$, deren Anzahl $C$ ist; die $2\nu - 7$ Constanten der Function $F(t)$; die reellen Gr\"ossen $b$, $c$, $e$, von denen man dreien beliebige Werthe geben kann, indem man f\"ur $t$ eine lineare Substitution mit reellen Coefficienten macht; die reellen und imagin\"aren Theile der Gr\"ossen $a$. Zu diesen willk\"urlichen Constanten kommen bei der Integration noch $10$ hinzu, n\"amlich, wenn $\eta = \frac{K_{2} - \alpha K_{1}}{K_{2} - \alpha' K_{1}}$ ist, die drei complexen Verh\"altnisse $\alpha:\beta:\gamma:\delta$, die f\"ur $6$ reelle Constanten zu z\"ahlen sind, ein (reeller oder rein imagin\"arer) Factor von $du$, und je eine additive reelle Constante in den Ausdr\"ucken f\"ur $x$, $y$, $z$. Diese Constanten m\"ussen aber noch Bedingungsgleichungen unterworfen werden, die erf\"ullt sein m\"ussen, wenn unsere Formeln wirklich eine Minimalfl\"ache darstellen sollen. Von diesen Bedingungsgleichungen sind $2A + B$ den Punkten $a$, $a'$, $b$ entsprechend, die aussagen, dass in den in der Umgebung dieser Punkte g\"ultigen Entwicklungen der L\"osungen der Differentialgleichung (15) keine Logarithmen auftreten (vergl.\ Anmerkung (4)), und $C + E$, die besagen, dass die zwischen den einzelnen Punkten $c$, $e$ gelegenen St\"ucke der Axe der reellen $t$ sich auf der Kugel mit dem Radius $1$ in $C + E$ B\"ogen gr\"osster Kreise abbilden. Sonach ist die Anzahl der in der L\"osung uebrig bleibenden unbestimmten Constanten $3C + 4E$.

Die Daten des Problems bestehen in den Coordinaten der Eckpunkte, und in den Winkeln, die die Richtungen der ins Unendliche verlaufenden Begrenzungslinien festlegen. Diese Daten sprechen sich in $3C + 4E$ Gleichungen aus, zu deren Erf\"ullung man eine ebenso grosse Zahl verf\"ugbarer Constanten hat.\textsuperscript{(*)})

\bigskip
\noindent\textbf{Beispiele.}

\subsection*{15.}

Die Begrenzung bestehe aus zwei unendlichen geraden Linien, die nicht in einer Ebene liegen. Ihre k\"urzeste Verbindungslinie habe die L\"ange $A$, und es sei $\alpha\pi$ der Winkel, welchen die Projection der Fl\"ache auf der rechtwinklig gegen jene Verbindungslinie gelegten Ebene ausf\"ullt.

Nimmt man die k\"urzeste Verbindungslinie zur $x$-Axe, so hat in jeder der beiden Begrenzungsgeraden $x$ einen constanten Werth. Ebenso ist $\varphi$ in jeder der beiden Begrenzungsgeraden constant. In unendlicher Entfernung ist die positive Normale f\"ur den einen Sector parallel der positiven, f\"ur den andern Sector parallel der negativen $x$-Axe. Die Begrenzung bildet sich auf der Kugel in zwei gr\"ossten Kreisen ab, die durch die Pole $\eta = 0$ und $\eta = \infty$ gehen und den Winkel $\alpha\pi$ einschliessen.

Hiernach hat man
\[x = \frac{iA}{2\alpha\pi} \log\eta,\]
\[s = \frac{iA}{2\alpha\pi\sqrt{\eta}} \left(\frac{1}{\eta} - \eta'\right),\]
\[s' = \frac{iA}{2\alpha\pi\sqrt{\eta'}} \left(\frac{1}{\eta'} - \eta\right),\]
folglich
\[x = -\frac{iA}{2\alpha\pi} \log(\eta\eta'),\]
\[y = \frac{A}{2\alpha\pi} \cdot \frac{1}{2i} \log\frac{s'}{s},\]
worin man die Gleichung der Schraubenfl\"ache erkennt.

\subsection*{16.}

Die Begrenzung bestehe aus drei geraden Linien, von denen zwei sich schneiden und die dritte zur Ebene der beiden ersten parallel l\"auft. Legt man den Anfangspunkt der Coordinaten in den Schnittpunkt der beiden ersten Geraden, die positive $x$-Axe in die negative Normale, so bildet jener Schnittpunkt auf der Kugel sich ab im Punkte $\eta = \infty$. Die Abbildung der beiden ersten Geraden sind gr\"osste Halbkreise, die von $\eta = \infty$ bis $\eta = 0$ laufen. Ihr Winkel sei $\alpha\pi$. Die Abbildung der dritten Linie ist der Bogen eines gr\"ossten Kreises, der von $\eta = 0$ ausgeht, an einer gewissen Stelle umkehrt und in sich selbst bis zum Punkte $\eta = 0$ zur\"uckl\"auft. Dieser Bogen bilde mit den beiden ersten gr\"ossten Halbkreisen die Winkel $-\beta\pi$ und $\gamma\pi$, so dass $\beta$ und $\gamma$ absolute Zahlen sind und $\beta + \gamma = \alpha$ sich ergiebt. Um die Abbildung auf der halben $t$-Ebene zu erhalten, setzen wir fest, dass $t = \infty$ sein soll f\"ur $\eta = \infty$, dass dem unendlichen Sector zwischen der ersten und dritten Linie $t = b$, dem unendlichen Sector zwischen der zweiten und dritten Linie $t = c$, dem Umkehrpunkte der Normalen auf der dritten Linie $t = a$ entsprechen soll. Dabei sind $a$, $b$, $c$ reell und $c > a > b$.

Diesen Bestimmungen entspricht $\eta = (t-b)^{\beta}(t-c)^{\gamma}$. Der Werth $a$ h\"angt von $b$ und $c$ ab. Man hat n\"amlich
\[\frac{d\log\eta}{dt} = \frac{\beta(t-c) + \gamma(t-b)}{(t-b)(t-c)}\]
und dieses muss f\"ur den Umkehrpunkt $=0$ sein, also $a = \frac{\beta c + \gamma b}{\beta + \gamma}$.

Man hat weiter nach Art.~12 und 13
\[\frac{du}{dt} = i\sqrt{(c-b)(\beta+\gamma)} \frac{(t-a)\, dt}{(t-b)(t-c)},\]
oder wenn man $c - b = -r$ annimmt
\[\frac{du}{dt} = i\sqrt{(c-b)(\beta+\gamma)} \frac{(t-a)\, dt}{(t-b)(t-c)}.\]

\newpage
Folglich
\[\text{(b)} \quad X = i\sqrt{\beta+\gamma} \int \frac{(t-a)\, dt}{\sqrt{(t-b)(t-c)}},\]
\[\frac{ds}{dt} = i\sqrt{\beta+\gamma} \cdot \frac{1}{\eta^{2}} \frac{(t-a)\, dt}{(t-b)(t-c)},\]
\[\frac{ds'}{dt} = i\sqrt{\beta+\gamma} \cdot \eta^{2} \frac{(t-a)\, dt}{(t-b)(t-c)},\]
\[s = \frac{i}{2} \int \frac{(t-b)^{\beta}(t-c)^{\gamma} - (t-b)^{-\beta}(t-c)^{-\gamma}}{(t-b)(t-c)}\, (t-a)\, dt,\]
\[\quad + \frac{i}{2} \int \frac{(t'-b)^{\beta}(t'-c)^{\gamma} - (t'-b)^{-\beta}(t'-c)^{-\gamma}}{(t'-b)(t'-c)}\, (t'-a)\, dt',\]
\[\frac{(t-b)^{\beta}(t-c)^{\gamma} + (t-b)^{-\beta}(t-c)^{-\gamma}}{(t-b)(t-c)}\, (t-a)\, dt.\]

\subsection*{17.}

Die Begrenzung bestehe aus drei einander kreuzenden geraden Linien, deren k\"urzeste Abst\"ande $A$, $B$, $C$ sein m\"ogen. Zwischen je zwei begrenzenden Linien erstreckt sich die Fl\"ache ins Unendliche. Es seien $\alpha\pi$, $\beta\pi$, $\gamma\pi$ die Winkel der Richtungen, in welchen die Grenzlinien des ersten, des zweiten, des dritten Sectors ins Unendliche verlaufen. Setzt man fest, dass f\"ur die drei Sectoren der Minimalfl\"ache im Unendlichen die Gr\"osse $t$ resp.\ $= 0$, $\infty$, $1$ sein soll, so erh\"alt man
\[\frac{du}{dt} = \frac{\sqrt{\varphi(t)}}{t(1-t)},\]
$\varphi(t)$ ist eine ganze Function zweiten Grades. Ihre Coefficienten bestimmen sich daraus, dass
\[\text{f\"ur } t = 0: \quad \frac{du}{d\log t} = \frac{A}{2\pi},\]
\[\text{f\"ur } t = \infty: \quad \frac{du}{d\log t} = \frac{B}{2\pi},\]
\[\text{f\"ur } t = 1: \quad \frac{du}{d\log(1-t)} = \frac{C}{2\pi}\]
sein muss.

Danach ergiebt sich
\[\varphi(t) = \frac{1}{4\pi^{2}} \left\{A^{2} + (B^{2} - A^{2} - C^{2})t + C^{2}t^{2}\right\}.\]
Je nachdem die Wurzeln der Gleichung $\varphi(t) = 0$ imagin\"ar oder reell sind, hat die Abbildung auf der Kugel einen Verzweigungspunkt im Innern oder zwei Umkehrpunkte der Normalen auf der Begrenzung.

\newpage
Die Functionen $k_{1} = \sqrt{\frac{dv}{d\eta}}$ und $k_{2} = \eta\sqrt{\frac{dv}{d\eta}}$ werden nur f\"ur die drei Sectoren unstetig, wenn man $\frac{dv}{dt} = \varphi(t)$ nimmt, und zwar ist die Unstetigkeit von $k_{1}$ der Art, dass
\[\text{f\"ur } t = 0: \quad k_{1} = t^{-\frac{\alpha}{2}},\]
\[\text{f\"ur } t = \infty: \quad k_{1} = t^{\frac{\beta}{2}},\]
\[\text{f\"ur } t = 1: \quad k_{1} = (1-t)^{\frac{\gamma}{2}}\]
ein\"andrig und verschieden von $0$ und $\infty$ wird. $k_{1}$ und $k_{2}$ sind particul\"are Integrale einer homogenen linearen Differentialgleichung zweiter Ordnung, die sich ergiebt, wenn man $\frac{1}{k} \frac{d^{2}k}{dv^{2}}$ aus seinen Unstetigkeiten als Function von $t$ darstellt und $t$ statt $v$ als unabh\"angige Variable in $\frac{d^{2}k}{dv^{2}}$ einf\"uhrt. Hat man das particul\"are Integral $k_{1}$ gefunden, so ergiebt sich $k_{2}$ aus der Differentialgleichung erster Ordnung
\[k_{2} \frac{dk_{1}}{dv} - k_{1} \frac{dk_{2}}{dv} = 1.\]
Das vollst\"andige Integral der homogenen linearen Differentialgleichung zweiter Ordnung werde mit
\[\text{(d)} \quad k = P \begin{pmatrix} 0 & \infty & 1 \\ -\frac{\alpha}{2} & +\frac{\beta}{2} & \frac{\gamma}{2} & t \\ \ast & \ast & \ast \end{pmatrix}\]
bezeichnet. Diese Function gen\"ugt wesentlich denselben Bedingungen, die in der Abhandlung ueber die Gauss'sche Reihe $F(\alpha, \beta, \gamma, x)$ als Definition der $P$-Function ausgesprochen sind\textsuperscript{*}). Sie weicht von der $P$-Function darin ab, dass die Summe der Exponenten $-1$ ist, nicht $+1$ wie bei $P$.

Man kann die Function $Q$ mit H\"ulfe einer Function $P$ und ihrer ersten Derivirten ausdr\"ucken. Zun\"achst ist n\"amlich
\[k_{1} = P \begin{pmatrix} 0 & \infty & 1 & \\ -\frac{\alpha}{2} & +\frac{\beta}{2} & \frac{\gamma}{2} & t \\ \frac{\alpha}{2} & -\frac{\beta}{2}-\gamma-1 & -\frac{\gamma}{2} \end{pmatrix},\]
\[Q = P \begin{pmatrix} 0 & \infty & 1 \\ 0 & -\alpha-\beta-\gamma+1 & 0 & t \\ \alpha & -\alpha+\beta-\gamma+1 & \gamma \end{pmatrix}.\]

\newpage
Setzt man nun
\[k_{1} = t^{-\frac{\alpha}{2}} (1-t)^{\frac{\gamma}{2}} \vartheta,\]
\[Q = P \begin{pmatrix} 0 & \infty & 1 \\ 0 & -\alpha-\beta-\gamma+1 & 0 & t \\ \alpha & -\alpha+\beta-\gamma+1 & \gamma \end{pmatrix},\]
so lassen sich die Constanten $a$, $b$, $c$ so bestimmen, dass
\[\text{(e)} \quad k_{1} = t^{-\frac{\alpha}{2}} (1-t)^{\frac{\gamma}{2}} \left((a+bt)\vartheta + ct(1-t) \frac{d\vartheta}{dt}\right)\]
wird. In der That hat man nur diesen Ausdruck in die Differentialgleichung (c) einzusetzen und die Differentialgleichung zweiter Ordnung f\"ur $\vartheta$ zu beachten, um zu der Gleichung zu gelangen
\[\varphi(t) = a^{2} t^{-\alpha} (1-t)^{-\gamma} (\vartheta, \vartheta)\, dt,\]
\[F(t) = a(a+\alpha)(1-t) + (a+b)(a+b-\gamma c)t\]
\[= (1-t)\left(1 - \frac{a+b+c}{a}t\right)\left(1 - \frac{a+b-c}{a}t\right).\]
Verm\"oge der Eigenschaften der Function $\vartheta$ kann man setzen
\[t^{-\alpha}(1-t)^{-\gamma}(\vartheta, \vartheta) = 1.\]
und folglich muss $F(t) = \varphi(t)$ sein. Hieraus ergeben sich drei Bedingungsgleichungen f\"ur $a$, $b$, $c$, die eine sehr einfache Form annehmen, wenn man
\[a+b+c = p, \quad a+b-c = q, \quad \frac{a+b+c}{a} = r\]
setzt. Die Bedingungsgleichungen lauten dann
\[pp - \alpha\alpha(p+q+r)^{2} = \frac{A^{2}}{2\pi^{2}},\]
\[qq - \beta\beta(p+q+r)^{2} = \frac{B^{2}}{2\pi^{2}},\]
\[rr - \gamma\gamma(p+q+r)^{2} = \frac{C^{2}}{2\pi^{2}}.\]
Mit H\"ulfe der Function
\[\varDelta = P \begin{pmatrix} 0 & \infty & 1 \\ \frac{\alpha}{2} & \frac{\beta}{2} & \frac{\gamma}{2} & t \\ -\frac{\alpha}{2} & -\frac{\beta}{2}+2 & -\frac{\gamma}{2}+2 \end{pmatrix},\]
deren Zweige $\varDelta_{1}$ und $\varDelta_{2}$ der Differentialgleichung gen\"ugen

\newpage
\noindent den Coordinatenaxen parallel laufen, ist $\alpha = \beta = \gamma = \frac{1}{2}$. Dann erh\"alt man
\[\text{(f)} \quad k = t^{\frac{1}{4}} \left((p+qt)^{2} + ct(1-t)^{2}\right).\]
Es w\"urde nicht schwer sein, die einzelnen Zweige der Function $k$ in der Form von bestimmten Integralen herzustellen. Der Weg dazu ist in Art.~7 der Abhandlung ueber die Function $P$ vorgezeichnet.

In dem besondern Falle, dass die drei begrenzenden geraden Linien den Coordinatenaxen parallel laufen, ist $\alpha = \beta = \gamma = \frac{1}{2}$. Dann erh\"alt man

Der Zweig $K_{1}$ dieser Function ist
\[= \left(\frac{c}{p}\right)^{\frac{1}{2}} \sqrt{t^{2} + (q-1)^{2}} \cdot \text{const.},\]
und daraus ergiebt sich
\[\begin{aligned}
k_{1} &= \sqrt{2} t^{\frac{1}{4}} (t-1)^{\frac{1}{4}} \sqrt{t^{2} + (q-1)^{2}} \left[p + qt + i\sqrt{c}\, t\sqrt{t^{2}+(q-1)^{2}}\right], \\
k_{2} &= \sqrt{2} t^{\frac{1}{4}} (t-1)^{\frac{1}{4}} \sqrt{t^{2} + (q-1)^{2}} \left[p + qt - i\sqrt{c}\, t\sqrt{t^{2}+(q-1)^{2}}\right].
\end{aligned}\]

Mit H\"ulfe dieser beiden Functionen lassen sich $dX$, $dY$, $dZ$ folgendermassen ausdr\"ucken
\[dX = -i\, k_{1}k_{2}\, \frac{du}{d\log\eta}\, dt,\]
\[dZ = -\frac{i}{2} (k_{1}^{2} + k_{2}^{2}) \frac{du}{d\log\eta}\, dt,\]
\[\text{(g)} \quad iY = -(p-q+r)\sqrt{t} + (p+q-r)\sqrt{t^{2}}\]
\[\quad - i(p+3q+3r)(p+q-r)\log\sqrt{t+1},\]
\[iZ = (p-q+r)^{2}(1-t)^{\frac{1}{2}} + (p+q-r)^{2}(1-t)^{-\frac{1}{2}}\]
\[\quad + i(3p+q+r)(-p+q+r)\log\sqrt{1+t^{2}}.\]

\newpage
Wenn $p$, $q$, $r$ reell sind, so geben die doppelten Coefficienten von $i$ in den drei Gr\"ossen rechts die rechtwinkligen Coordinaten eines Punktes der Fl\"ache.

\subsection*{18.}

Die Begrenzung bestehe aus vier sich schneidenden geraden Linien, die man erh\"alt\textsuperscript{**}) wenn von den Kanten eines beliebigen Tetraeders zwei nicht zusammenstossende weggelassen werden. Die Abbildung auf der Kugeloberfl\"ache ist ein sph\"arisches Viereck, dessen Winkel $\alpha\pi$, $\beta\pi$, $\gamma\pi$, $\delta\pi$ sein m\"ogen. Es ergiebt sich
\[\frac{du}{dt} = \frac{G\, dt}{\sqrt{(t-a)(t-b)(t-c)(t-d)}} = \frac{G\, dt}{\sqrt{\varDelta(t)}},\]
wenn die reellen Werthe $t = a$, $b$, $c$, $d$ die Punkte der $t$-Ebene bezeichnen, in welchen sich die Eckpunkte des Vierecks abbilden.

Soll die in \S.~14 entwickelte Methode zur Bestimmung von $\eta$ angewandt werden, so hat man hier speciell $\varphi(t) = 1$, $\chi(t) = \varDelta(t)$, folglich $\nu = 4$ und
\[\frac{dv}{dt} = \frac{1}{\sqrt{\varDelta(t)}}.\]
Die Functionen $k_{1}$ und $k_{2}$ gen\"ugen der Differentialgleichung
\[k_{2} \frac{dk_{1}}{dv} - k_{1} \frac{dk_{2}}{dv} = 1\]
und sind particul\"are Integrale der Differentialgleichung zweiter Ordnung
\[\frac{1}{k} \frac{d^{2}k}{dv^{2}} = \frac{(\alpha\alpha-1)\varDelta'(a)}{4(t-a)} + \frac{(\beta\beta-1)\varDelta'(b)}{4(t-b)}\]
\[+ \frac{(\gamma\gamma-1)\varDelta'(c)}{4(t-c)} + \frac{(\delta\delta-1)\varDelta'(d)}{4(t-d)}.\]
Die Function $F(t)$ des \S.~14 ist hier vom zweiten Grade, aber die Coefficienten von $t^{2}$ und von $t$ sind gleich Null, also $h$ eine Constante.

In der letzten Gleichung hat man auf der linken Seite $t$ als unabh\"angige Variable einzuf\"uhren und erh\"alt
\[\frac{1}{k} \frac{d^{2}k}{dt^{2}} = \frac{(\alpha\alpha-1)\varDelta'(a)}{t-a} + \frac{(\beta\beta-1)\varDelta'(b)}{t-b}\]
\[+ \frac{(\gamma\gamma-1)\varDelta'(c)}{t-c} + \frac{(\delta\delta-1)\varDelta'(d)}{t-d}\]
als die Differentialgleichung zweiter Ordnung, welcher $k$ Gen\"uge leisten muss.

Sind $X$, $Y$, $Z$ als Functionen von $t$ wirklich ausgedr\"uckt, so treten in der L\"osung noch $16$ unbestimmte reelle Constanten auf, n\"amlich die vier Gr\"ossen $a$, $b$, $c$, $d$, von denen wie oben drei beliebig angenommen werden k\"onnen, die vier Gr\"ossen $\alpha$, $\beta$, $\gamma$, $\delta$, die Gr\"osse $h$,

\newpage
\noindent ferner $6$ reelle Constanten in dem Ausdrucke f\"ur $\eta$, ein constanter Factor in $du$ und je eine additive Constante in $x$, $y$, $z$. Zur Bestimmung dieser $16$ Gr\"ossen sind $16$ Bedingungsgleichungen vorhanden, n\"amlich $4$ Gleichungen, welche ausdr\"ucken, dass die vier Begrenzungslinien in der Ebene der $\eta$ sich auf der Kugel in gr\"ossten Kreisen abbilden, und $12$ Gleichungen, welche aussagen, dass $x$, $y$, $z$ in den $4$ Eckpunkten gegebene Werthe haben.

In dem speciellen Falle eines regul\"aren Tetraeders ist die Abbildung auf der Kugel ein regelm\"assiges Viereck, in welchem jeder Winkel $= \frac{2}{3}\pi$. Die Diagonalen halbiren sich und stehen rechtwinklig aufeinander. Die den Eckpunkten diametral gegen\"uberliegenden Punkte der Kugeloberfl\"ache sind die Ecken eines congruenten Vierecks. Zwischen beiden liegen vier dem urspr\"unglichen ebenfalls congruente Vierecke, die je zwei Eckpunkte mit dem urspr\"unglichen, zwei mit dem gegen\"uberliegenden gemein haben. Diese sechs Vierecke f\"ullen die Kugeloberfl\"ache einfach aus. Es wird also $\frac{du}{d\log\eta}$ eine algebraische Function von $\eta$ sein.

Man kann die gesuchte Minimalfl\"ache ueber ihre urspr\"ungliche Begrenzung dadurch stetig fortsetzen, dass man sie um jede ihrer Grenzlinien als Drehungsaxe um $180^{\circ}$ dreht. L\"angs einer solchen Grenzlinie haben dann die urspr\"ungliche Fl\"ache und die Fortsetzung gemeinschaftliche Normalen. Wiederholt man die Construction an den neuen Fl\"achentheilen, so l\"asst sich die urspr\"ungliche Fl\"ache beliebig weit fortsetzen. Welche Fortsetzung man aber auch betrachte, immer bildet sie sich auf der Kugel in einem der sechs congruenten Vierecke ab. Und zwar haben die Abbildungen von zwei Fl\"achentheilen eine Seite gemein oder sie liegen einander gegen\"uber, je nachdem die Fl\"achentheile selbst in einer Grenzlinie an einander stossen oder an gegen\"uberliegenden Grenzlinien eines mittleren Fl\"achentheils gelegen sind. In dem letzteren Falle k\"onnen die betreffenden Fl\"achentheile durch parallele Verschiebung zur Deckung gebracht werden. Daher muss $\left(\frac{du}{d\log\eta}\right)^{2}$ unver\"andert bleiben, wenn $\eta$ mit $\frac{1}{\eta'}$ vertauscht wird.

Legt man den Pol ($\eta = 0$) in den Mittelpunkt eines Vierecks, den Anfangsmeridian durch die Mitte einer Seite, so ist f\"ur die Eckpunkte dieses Vierecks
\[\eta = \tg\frac{\pi}{8} \cdot e^{\frac{\pi i}{4}}, \quad \tg\frac{\pi}{8} \cdot e^{\frac{3\pi i}{4}}, \quad \tg\frac{\pi}{8} \cdot e^{-\frac{3\pi i}{4}}, \quad \tg\frac{\pi}{8} \cdot e^{-\frac{\pi i}{4}},\]
und f\"ur die Eckpunkte des gegen\"uberliegenden Vierecks die entgegengesetzten Werthe.

\newpage
Punkte, denen entgegengesetzte Werthe von $\eta$ angeh\"oren, haben dieselbe $X$-Coordinate. Es muss also $\left(\frac{du}{d\log\eta}\right)^{2}$ bei der Vertauschung von $\eta$ mit $-\eta$ unver\"andert bleiben. Hiernach erh\"alt man
\[\left(\frac{du}{d\log\eta}\right)^{2} = C_{1}^{2} \frac{(\eta^{4}+\eta'^{4}+1)}{(\eta^{4}\eta'^{4}+\eta^{4}+\eta'^{4})}.\]
Die Constante $C_{1}^{2}$ muss reell sein, damit $du^{2}$ in der Begrenzung reelle Werthe besitze.

Zu demselben Resultate gelangt man auf dem folgenden Wege. Die Substitution
\[\xi^{2} = \frac{\eta^{4}+\eta'^{4}+1}{\eta^{4}\eta'^{4}+\eta^{4}+\eta'^{4}}\]
liefert auf der $\xi$-Ebene eine Abbildung, die von einer geschlossenen ueberall stetig gekr\"ummten Linie begrenzt wird. Die Rechnung zeigt, dass $d\log\xi$ in der Begrenzung rein imagin\"ar ist. Folglich ist die Abbildung der Begrenzung in der $\xi$-Ebene ein Kreis um den Mittelpunkt $\xi = 0$. Der Radius dieses Kreises ist $=1$. Den Eckpunkten
\[\eta = \tg\frac{\pi}{8} \cdot e^{\pm\frac{\pi i}{4}}\]
entspricht $\xi = +1$, den Eckpunkten
\[\eta = \tg\frac{\pi}{8} \cdot e^{\pm\frac{3\pi i}{4}}\]
entspricht $\xi = i$. Geht man an irgend einer dieser vier Stellen durch das Innere der Minimalfl\"ache von einer Grenzlinie zur folgenden, so \"andert sich dabei der Arcus von $d\xi$ um $\pi$. Daher kann man, wie in \S.~13, auch hier setzen
\[\frac{du}{dt} = \frac{C_{2}}{\sqrt{(t-a)(t-b)(t-c)}},\]
und es muss $C_{2}$ rein imagin\"ar sein, damit $du^{2}$ in der Begrenzung reell ausfalle. Es findet sich $C_{2} = 3\sqrt[4]{C_{1}^{2}} \cdot i$.

Dieser Ausdruck stimmt mit dem vorher aufgestellten f\"ur $\left(\frac{du}{d\log\eta}\right)^{2}$.

Zur weitern Vereinfachung nehme man
\[\eta^{2} = \frac{\sqrt{t^{2}+1} + \sqrt{2}\, t}{\sqrt{t^{2}+1} - \sqrt{2}\, t}\]
und beachte, dass
\[\left(\frac{du}{d\log\eta}\right)^{2} = 2C_{1}^{2} \frac{dt}{d\log\eta}.\]
Dann ergiebt eine sehr einfache Rechnung

\newpage
\begin{equation}\tag{h}
\frac{1}{2} \left(\frac{du}{d\log\eta}\right)^{2} = C \int \frac{dt}{\sqrt{1-t^{4}}},
\end{equation}
\[\frac{1}{2} \left(\frac{du}{d\log\eta}\right)^{2} = C \int \frac{dt}{\sqrt{(1-\varrho t^{2})(1-\varrho^{2}t^{2})}},\]
wenn $\varrho = -\frac{1}{2}(1-i\sqrt{3})$ eine dritte Wurzel der Einheit bezeichnet.

Die reelle Constante $C = \sqrt[3]{6} C_{1}^{2}$ bestimmt sich aus der gegebenen L\"ange der Tetraederkanten.

\subsection*{19.}

Endlich soll noch die Aufgabe der Minimalfl\"ache f\"ur den Fall behandelt werden, dass die Begrenzung aus zwei beliebigen Kreisen besteht, die in parallelen Ebenen liegen. Dann kennt man die Richtung der Normalen in der Begrenzung nicht. Daher l\"asst sich diese auch nicht auf der Kugel abbilden. Man gelangt aber zur L\"osung der Aufgabe durch die Annahme, dass alle zu den Ebenen der Grenzkreise parallel gelegten ebenen Schnitte Kreise seien, und es wird sich zeigen, dass unter dieser Annahme der Minimalbedingung Gen\"uge geleistet werden kann.

Legt man die $x$-Axe rechtwinklig gegen die Ebenen der Grenzkreise, so ist die Gleichung der Schnittcurve in einer parallelen Ebene
\[y^{2} + z^{2} + 2\alpha y + 2\beta z + \gamma = 0,\]
und $\alpha$, $\beta$, $\gamma$ sind als Functionen von $x$ zu bestimmen. Zur Abk\"urzung werde
\[F = y^{2} + z^{2} + 2\alpha y + 2\beta z + \gamma\]
gesetzt, so dass
\[\cos r = w \frac{\partial F}{\partial x}, \quad \sin r \cos\varphi = w \frac{\partial F}{\partial y}, \quad \sin r \sin\varphi = w \frac{\partial F}{\partial z}\]
ist. Dann l\"asst sich die Bedingung des Minimum in die Form bringen
\[\frac{\partial}{\partial x} \left(\frac{1}{w}\right) + \frac{\partial}{\partial y} \left(\frac{\sin r \cos\varphi}{w}\right) + \frac{\partial}{\partial z} \left(\frac{\sin r \sin\varphi}{w}\right) = 0\]
oder nach Ausf\"uhrung der Differentiation
\[\frac{\partial^{2}F}{\partial y^{2}} + \frac{\partial^{2}F}{\partial z^{2}} + \frac{w'}{w} \frac{\partial F}{\partial x} + 4w^{2}(F + \alpha^{2} + \beta^{2} - \gamma) = 0.\]
Schreibt man $\alpha^{2} + \beta^{2} - \gamma = -q$ und beachtet, dass $F = 0$ ist, so geht die letzte Gleichung ueber in

\newpage
\noindent $\frac{w'}{w} + 2q + \text{const.} = 0$

und giebt nach einmaliger Integration
\[\left|\frac{dx}{w}\right| + 2\int q\, dx + \text{const.} = 0.\]
Die Integrationsconstante ist von $x$ unabh\"angig. Nimmt man andererseits $\frac{dx}{w}$ unabh\"angig von $y$ und $z$, so muss die Integrationsconstante eine lineare Function von $y$ und $z$ sein, weil $2\int q\, dx$ eine solche ist. Man hat also
\[\frac{1}{w} + 2\int q\, dx + 2ay + 2bz + \text{const.} = 0.\]
Vergleicht man damit das Resultat der directen Differentiation von $F$, n\"amlich
\[\frac{\partial F}{\partial x} = 2y \frac{d\alpha}{dx} + 2z \frac{d\beta}{dx} + \frac{d\gamma}{dx},\]
so ergiebt sich
\[\frac{d\alpha}{dx} = -aw, \quad \frac{d\beta}{dx} = -bw,\]
und wenn man $\int q\, dx = m$ setzt:
\[\alpha = -am + d, \quad \beta = -bm + e.\]
Hiernach hat man
\[q = -2a\varrho y - 2b\varrho z + f(m),\]
\[\frac{d^{2}q}{dx^{2}} = -2a\varrho \frac{dy}{dx} - 2b\varrho \frac{dz}{dx} + f'(m)\frac{dm}{dx} + f''(m)\left(\frac{dm}{dx}\right)^{2}\]
und diese Ausdr\"ucke sind in die Gleichung (1) einzuf\"uhren. Nach geh\"origer Hebung erh\"alt man
\[\frac{d^{2}q}{dx^{2}} + 2q \frac{dm}{dx} + 2q = 0\]
eine Gleichung, die sich weiter vereinfacht, wenn man beachtet, dass
\[f(m) = (a^{2} + b^{2})m^{2} - 2(ad + be)m + d^{2} + e^{2}.\]
Nimmt man hieraus $\frac{dq}{dx}$ und $\frac{d^{2}q}{dx^{2}}$, so geht die Differentialgleichung, welche die Bedingung des Minimum ausdr\"uckt, ueber in folgende:
\[\frac{d^{2}m}{dx^{2}} + 2m + 2(a^{2} + b^{2})m^{2} = 0.\]
Zur Ausf\"uhrung der Integration setze man $\frac{dm}{dx} = p$ und betrachte

\newpage
\noindent $q$ als unabh\"angige Variable. Dadurch erh\"alt man f\"ur $p^{2}$ als Function von $q$ eine lineare Differentialgleichung erster Ordnung, n\"amlich
\[\frac{d(p^{2})}{dq} - p^{2} + 2q + 2(a^{2} + b^{2})q^{2} = 0\]
oder
\[\frac{d(p^{2})}{dq} - p^{2} = -2q - 2(a^{2} + b^{2})q^{2}.\]
Das Integral lautet
\[p^{2} = e^{q} \left\{ \text{const.} + \int e^{-q} \left[2q + 2(a^{2} + b^{2})q^{2}\right] dq \right\}.\]
Darin ist f\"ur $p$ wieder $\frac{dm}{dx}$ zu setzen, wodurch man erh\"alt
\[\frac{dx}{dm} = \frac{\sqrt{2q + 2cq^{2} - (a^{2} + b^{2})2q^{3}}}{1},\]
also ergiebt sich
\[\text{(o)} \quad m = \int \frac{dq}{\sqrt{2q + 2cq^{2} - (a^{2} + b^{2})2q^{3}}},\]
\[2y = am - d + \sqrt{-q}\, \cos\psi,\]
\[2z = bm - e + \sqrt{-q}\, \sin\psi.\]
Man hat demnach $x$, $y$, $z$ als Functionen von zwei reellen Variabeln $q$ und $\psi$ ausgedr\"uckt. Die Ausdr\"ucke sind, abgesehen von algebraischen Gliedern, elliptische Integrale mit der obern Grenze $q$. Nach der oben entwickelten allgemeinen Methode h\"atte man $x$, $y$, $z$ erhalten als Summen von zwei conjugirten Functionen zweier conjugirter complexer Variabeln. Danach liegt die Vermuthung nahe, dass diese complexen Ausdr\"ucke mit H\"ulfe der Additionstheoreme der elliptischen Functionen sich je in einen einzigen Integralausdruck mit der Variabeln $q$ zusammenziehen lassen.

Und dies ist leicht zu best\"atigen. Man hat n\"amlich aus den Formeln f\"ur die Richtungscoordinaten $r$ und $\varphi$ der Normalen
\[\sqrt{\frac{\partial F}{\partial y} + \frac{\partial F}{\partial z} i} = \sqrt{y + zi + \alpha + \beta i} = \sqrt{\eta}\, e^{\frac{r}{2}i}.\]
Verbindet man damit die Definitionsgleichung von $q$, n\"amlich:
\[\sqrt{(y + zi + \alpha + \beta i)(y - zi + \alpha - \beta i)} = -q,\]
so ergiebt sich

\newpage
\noindent $(y + zi) + (\alpha + \beta i) = (-q)^{\frac{1}{2}} \eta^{\frac{1}{2}} e^{\frac{r}{2}i},$
$(y - zi) + (\alpha - \beta i) = (-q)^{\frac{1}{2}} \eta^{-\frac{1}{2}} e^{-\frac{r}{2}i}.$

Ferner hat man
\[\cot\frac{r}{2} \frac{dx}{\sqrt{-(dy^{2} + dz^{2})}} = \frac{1}{\sqrt{-q}} \frac{dx}{d\log\eta}\]
oder
\[\sqrt{nn'} = \frac{1}{\sqrt{-q}} \frac{dx}{d\log\eta},\]
\[2\sqrt{-2}\left\{p - 2aq(y + \alpha) - 2bq(z + \beta)\right\}\]
\[= \sqrt{\eta}\left(\cos\frac{r}{2} - \sin\frac{r}{2}\right) - \frac{1}{\sqrt{\eta}}\left(\sin\frac{r}{2} + \cos\frac{r}{2}\right)\]
\[\cdot \sqrt{-2}\left\{p - 2a^{2}(y + \alpha) - 2b^{2}(z + \beta)\right\}.\]
Auf der rechten Seite sind f\"ur $y + \alpha$ und $z + \beta$ die eben gefundenen Ausdr\"ucke in $\eta$ und $\eta'$ einzuf\"uhren. Dadurch geht die Gleichung ueber in folgende:
\[\sqrt{-q}\, \frac{d\log\eta}{dx} = \sqrt{2c + (a+bi)\eta^{-1} - (a-bi)\eta}.\]
Quadrirt man beide Seiten dieser Gleichung und setzt f\"ur $\frac{dx}{d\log\eta}$ seinen Werth aus (n), so ergiebt sich nach geh\"origer Reduction
\[\text{(p)} \quad (1 + \eta\eta')^{2} \left\{2c - (a+bi)\frac{\eta'}{\eta} - (a-bi)\frac{\eta}{\eta'}\right\}\]
\[= 8c - 2(a+bi)(\eta' - \eta^{-1}) - 2(a-bi)(\eta - \eta'^{-1}).\]
Die so gefundene Gleichung, welche den Zusammenhang von $q$, $\eta$, $\eta'$ angiebt, kann man als Integral einer Differentialgleichung f\"ur $\eta$ und $\eta'$ ansehen und $q$ als Integrationsconstante auffassen. Die Differentialgleichung ergiebt sich durch unmittelbare Differentiation in folgender Form
\[0 = \frac{d\eta}{\eta} \sqrt{2c + (a+bi)\eta^{-1} - (a-bi)\eta}\]
\[- \frac{d\eta'}{\eta'} \sqrt{2c + (a-bi)\eta'^{-1} - (a+bi)\eta'}.\]
Mit H\"ulfe der primitiven Gleichung (p) lassen sich aber die Factoren von $\frac{d\eta}{\eta}$ und $\frac{d\eta'}{\eta'}$ anders ausdr\"ucken. Man braucht nur die linke Seite

\newpage
\noindent von (p) in zweifacher Weise zu einem vollst\"andigen Quadrat zu erg\"anzen, indem man das fehlende doppelte Product das eine mal positiv, das andere mal negativ hinzuf\"ugt. Dadurch erh\"alt man
\[= +2\sqrt{[2c + (a+bi)\eta^{-1}][2c - (a-bi)\eta]}\]
\[-2\sqrt{[2c - (a+bi)\eta^{-1}][2c + (a-bi)\eta]}.\]
Nimmt man die Quadratwurzeln mit gleichen Vorzeichen, so geht die Differentialgleichung ueber in
\[0 = \frac{d\eta}{2\eta\sqrt{2c + (a+bi)\eta^{-1} - (a-bi)\eta}}\]
\[\text{(q)} \quad - \frac{d\eta'}{2\eta'\sqrt{2c + (a-bi)\eta'^{-1} - (a+bi)\eta'}}.\]
Ihr Integral in algebraischer Form ist in der Gleichung (p) ausgesprochen oder, was auf dasselbe hinauskommt, in den beiden Gleichungen
\[\sqrt{-q}\, (1 + \eta\eta') = \sqrt{[2c + (a+bi)\eta + (a-bi)\eta']}\]
\[\text{(r)} \quad + \sqrt{[2c - (a+bi)\eta - (a-bi)\eta']},\]
\[\sqrt{-q}\, i[(a+bi)\eta - (a-bi)\eta'] = \sqrt{[2c(a+bi) + 2c\eta' - (a-bi)\eta'^{2}]}\]
\[- \sqrt{[2c(a-bi) + 2c\eta - (a+bi)\eta^{2}]}.\]
In transcendenter Form lautet das Integral
\[\text{const.} = \int \frac{d\eta}{2\eta\sqrt{(a+bi) + 2c\eta - (a-bi)\eta^{2}}}\]
\[+ \int \frac{d\eta'}{2\eta'\sqrt{(a-bi) + 2c\eta' - (a+bi)\eta'^{2}}},\]
und die Integrationsconstante l\"asst sich ausdr\"ucken
\[\text{const.} = \int \frac{dq}{2\sqrt{q[1 + 2cq - (a^{2} + b^{2})q^{2}]}},\]
was aus der Gleichung (r) leicht hervorgeht, wenn man $\eta$ oder $\eta'$ constant und zwar $= 0$ nimmt. Man erkennt darin das Additionstheorem der elliptischen Integrale erster Gattung.

\newpage
%==============================================================================
% ANMERKUNGEN (Notes)
%==============================================================================
\section*{Anmerkungen.}
\addcontentsline{toc}{section}{Anmerkungen.}

Die erste Ausgabe der Abhandlung ueber Fl\"achen kleinsten Inhalts durch Hattendorff in den Abhandlungen der Gesellschaft der Wissenschaften zu G\"ottingen vom Jahre 1867 war eingeleitet durch eine historische Betrachtung, die bereits in der ersten Auflage dieser Gesammtausgabe als nicht von Riemann herr\"uhrend in Uebereinstimmung mit dem leider fr\"uh verstorbenen Hattendorff unterdr\"uckt wurde. Sie ist auch in der zweiten Auflage nicht mit aufgenommen. Zu erw\"ahnen ist aber, dass ungef\"ahr gleichzeitig mit Riemann Weierstrass Untersuchungen ueber die Fl\"achen vom kleinsten Inhalt angestellt hat, deren Resultate in den Monatsberichten der Berliner Akademie vom October und December 1866 zuerst ver\"offentlicht sind. Die Weierstrass'schen Arbeiten haben den Anstoss gegeben zu den weitergehenden Untersuchungen von H.~A.~Schwarz, von denen die erste Mittheilung im Monatsberichte der Berliner Akademie vom April 1865 gemacht ist. Die ausf\"uhrliche Ver\"offentlichung der gekr\"onten Preisschrift ``Bestimmung einer speciellen Minimalfl\"ache'' erschien im Jahre 1871. Darin ist neben Anderem das in Art.~18 der Riemann'schen Abhandlung behandelte Problem der durch ein regelm\"assiges r\"aumliches Viereck begrenzten Minimalfl\"ache bis zur wirklichen Aufstellung einer Gleichung zwischen den Coordinaten $x$, $y$, $z$ der Fl\"ache durchgef\"uhrt. Zu erw\"ahnen ist noch die an Riemann direct ankn\"upfende, im Jahre 1867 von der philosophischen Facult\"at zu G\"ottingen gekr\"onte Preisschrift von Arthur Schondorff ``Ueber die Minimalfl\"ache, deren Begrenzung von einem doppeltgleichschenkligen r\"aumlichen Viereck gebildet wird''.

Die Riemann'sche Abhandlung ist hier, mit wenigen nothwendigen Aenderungen, in der letzten Hattendorff'schen Bearbeitung abgedruckt. Einige Erl\"auterungen und Zus\"atze gebe ich in den nachstehenden Anmerkungen, wobei ich Mittheilungen und Winke von H.~A.~Schwarz mit Dank benutzt habe.

\medskip
\noindent\textbf{(1)} (Zu Seite 307.) Es ist hierbei zu beachten, dass nach (3) und (4)
\[\frac{dX}{d\eta} = \frac{1}{\eta'} \frac{ds}{d\eta}, \quad \frac{dX}{d\eta'} = \frac{1}{\eta} \frac{ds}{d\eta'}\]
Functionen von $\eta$ allein sind, und dass also auch $Y$, $Z$ als Functionen von $\eta$ allein betrachtet werden k\"onnen. $Y'$, $Z'$ sind dann die conjugirten Functionen, die nur von $\eta'$ abh\"angen.

\medskip
\noindent\textbf{(2)} (Zu Seite 311.) F\"ur unendlich kleine Werthe von $\eta$ (also auch von $r$) ergiebt sich aus (1) und (2)
\[\frac{dx}{dy} = -\sin r \cos\varphi, \quad \frac{dx}{dz} = -\sin r \sin\varphi,\]
\[\frac{dx}{dy} = -\sin r \sin\varphi, \quad \frac{dx}{dz} = \sin r \cos\varphi.\]

\end{document}
